\section{AIME I 2022}
        \begin{enumerate}
        	\item (\textit{AIME I 2022}) Quadratic polynomials $P(x)$ and $Q(x)$ have leading coefficients $2$ and $-2,$ respectively. The graphs of both polynomials pass through the two points $(16,54)$ and $(20,53).$ Find $P(0) + Q(0).$


 
 \textbf{Answer: }116



\textbf{Solution 1 (Linear Polynomials)}

Let $R(x)=P(x)+Q(x).$ Since the $x^2$-terms of $P(x)$ and $Q(x)$ cancel, we conclude that $R(x)$ is a linear polynomial.

 Note that
 \begin{alignat*}{8} R(16) &= P(16)+Q(16) &&= 54+54 &&= 108, \\ R(20) &= P(20)+Q(20) &&= 53+53 &&= 106, \end{alignat*}so the slope of $R(x)$ is $\frac{106-108}{20-16}=-\frac12.$

 It follows that the equation of $R(x)$ is 
 \[R(x)=-\frac12x+c\] for some constant $c,$ and we wish to find $R(0)=c.$

 We substitute $x=20$ into this equation to get $106=-\frac12\cdot20+c,$ from which $c=\boxed{116}.$

 





\textbf{Solution 2 (Quadratic Polynomials)}

Let
 \begin{alignat*}{8} P(x) &= &2x\textasciicircum{}2 + ax + b, \\ Q(x) &= &\hspace{1mm}-2x\textasciicircum{}2 + cx + d, \end{alignat*}for some constants $a,b,c$ and $d.$ 

 We are given that
 \begin{alignat*}{8} P(16) &= &512 + 16a + b &= 54, \hspace{20mm}&&(1) \\ Q(16) &= &\hspace{1mm}-512 + 16c + d &= 54, &&(2) \\ P(20) &= &800 + 20a + b &= 53,  &&(3) \\ Q(20) &= &\hspace{1mm}-800 + 20c + d &= 53, &&(4) \end{alignat*}and we wish to find 
 \[P(0)+Q(0)=b+d.\]We need to cancel $a$ and $c.$ Since $\operatorname{lcm}(16,20)=80,$ we subtract $4\cdot[(3)+(4)]$ from $5\cdot[(1)+(2)]$ to get 
 \[b+d=5\cdot(54+54)-4\cdot(53+53)=\boxed{116}.\]~MRENTHUSIASM

 





\textbf{Solution 3 (Similar to Solution 2)}

Like Solution 2, we can begin by setting $P$ and $Q$ to the quadratic above, giving us
 \begin{alignat*}{8} P(16) &= &512 + 16a + b &= 54, \hspace{20mm}&&(1) \\ Q(16) &= &\hspace{1mm}-512 + 16c + d &= 54, &&(2) \\ P(20) &= &800 + 20a + b &= 53,  &&(3) \\ Q(20) &= &\hspace{1mm}-800 + 20c + d &= 53, &&(4) \end{alignat*}We can first add $(1)$ and $(2)$ to obtain $16(a-c) + (b+d) = 108.$ 

 Next, we can add $(3)$ and $(4)$ to obtain $20(a-c) + (b+d) = 106.$ By subtracting these two equations, we find that $4(a-c) = -2,$ so substituting this into equation $[(1) + (2)],$ we know that $4 \cdot (-2) + (b+d) = 108,$ so therefore $b+d = \boxed{116}.$

 





\textbf{Solution 4 (Brute Force)}

Let
 \begin{alignat*}{8} P(x) &= &2x\textasciicircum{}2 + ax + b, \\ Q(x) &= &\hspace{1mm}-2x\textasciicircum{}2 + cx + d, \end{alignat*}By substituting $(16, 54)$ and $(20, 53)$ into these equations, we can get:
 \begin{align*} 2(16)^2 + 16a + b &= 54, \\ 2(20)^2 + 20a + b &= 53. \end{align*}
Hence, $a = -72.25$ and $b = 698.$

 Similarly,
 \begin{align*} -2(16)^2 + 16c + d &= 54, \\ -2(20)^2 + 20c + d &= 53. \end{align*}
Hence, $c = 71.75$ and $d = -582.$

 Notice that $b = P(0)$ and $d = Q(0).$Therefore 
 \[P(0) + Q(0) = 698 + (-582) = \boxed{116}.\]~Littlemouse

 





\textbf{Solution 5}

Add the equations of the polynomials $y=2x^2+ax+b$ and $y=-2x^2+cx+d$ to get $2y=(a+c)x+(b+d)$. This equation must also pass through the two points $(16,54)$ and $(20,53)$.

 Let $m=a+c$ and $n=b+d$. We then have two equations:
 \begin{align*} 108&=16m+n, \\ 106&=20m+n. \end{align*}
We are trying to solve for $n=P(0)$. Using elimination:
 \begin{align*} 540&=80m+5n, \\ 424&=80m+4n. \end{align*}
Subtracting both equations, we find that $n=\boxed{116}$.

 




	\item (\textit{AIME I 2022}) Find the three-digit positive integer $\underline{a}\,\underline{b}\,\underline{c}$ whose representation in base nine is $\underline{b}\,\underline{c}\,\underline{a}_{\,\text{nine}},$ where $a,$ $b,$ and $c$ are (not necessarily distinct) digits.


 
 \textbf{Answer: }227



\textbf{Solution 1}

We are given that 
 \[100a + 10b + c = 81b + 9c + a,\] which rearranges to 
 \[99a = 71b + 8c.\]Taking both sides modulo $71,$ we have
 \begin{align*} 28a &\equiv 8c \pmod{71} \\ 7a &\equiv 2c \pmod{71}. \end{align*}
The only solution occurs at $(a,c)=(2,7),$ from which $b=2.$

 Therefore, the requested three-digit positive integer is $\underline{a}\,\underline{b}\,\underline{c}=\boxed{227}.$

 





\textbf{Solution 2}

As shown in Solution 1, we get $99a = 71b+8c$.

 Note that $99$ and $71$ are large numbers comparatively to $8$, so we hypothesize that $a$ and $b$ are equal and $8c$ fills the gap between them. The difference between $99$ and $71$ is $28$, which is a multiple of $4$. So, if we multiply this by $2$, it will be a multiple of $8$ and thus the gap can be filled. Therefore, the only solution is $(a,b,c)=(2,2,7)$, and the answer is $\underline{a}\,\underline{b}\,\underline{c}=\boxed{227}$.

 





\textbf{Solution 2a}

A little bit more motivation: taking mod $8$ on both sides, we find that $7a\equiv7b\pmod8\implies a\equiv b\pmod8$, so either $a=b$ or they differ by a multiple of $8$ (ok ok technically $a=b$ is a subcase of the other one, but shhhhh for the sake of clarity). $99$ and $71$ differ by $28$, so a difference of merely $8$ in $a$ and $b$ accounts for a huge difference on the scale of $28\cdot8=224$ (which is obviously unfillable by the mere $8c$), so $a=b$ is very reasonable. 

 





\textbf{Solution 3}

As shown in Solution 1, we get $99a = 71b+8c.$

 We list a few multiples of $99$ out: 
 \[99,198,297,396.\]Of course, $99$ can't be made of just $8$'s. If we use one $71$, we get a remainder of $28$, which can't be made of $8$'s either. So $99$ doesn't work. $198$ can't be made up of just $8$'s. If we use one $71$, we get a remainder of $127$, which can't be made of $8$'s. If we use two $71$'s, we get a remainder of $56$, which can be made of $8$'s. Therefore we get $99\cdot2=71\cdot2+8\cdot7$ so $a=2,b=2,$ and $c=7$. Plugging this back into the original problem shows that this answer is indeed correct. Therefore, $\underline{a}\,\underline{b}\,\underline{c}=\boxed{227}.$

 





\textbf{Solution 4}

As shown in Solution 1, we get $99a = 71b+8c$.

 We can see that $99$ is $28$ larger than $71$, and we have an $8c$. We can clearly see that $56$ is a multiple of $8$, and any larger than $56$ would result in $c$ being larger than $9$. Therefore, our only solution is $a = 2, b = 2, c = 7$. Our answer is $\underline{a}\,\underline{b}\,\underline{c}=\boxed{227}$.

 





\textbf{Solution 5}

As shown in Solution 1, we get $99a = 71b+8c,$ which rearranges to 
 \[99(a – b) = 8c – 28 b = 4(2c – 7b) \le 4(2\cdot 9 - 0 ) = 72.\]So $a=b, 2c = 7b \implies c=7, b=2,a=2.$

 \textbf{vladimir.shelomovskii@gmail.com, vvsss}

 





\textbf{Solution 6}

As shown in Solution 1, we have that $99a = 71b + 8c$.

 Note that by the divisibility rule for $9$, we have $a+b+c \equiv a \pmod{9}$.  Since $b$ and $c$ are base-$9$ digits, we can say that $b+c = 0$ or $b+c=9$.  The former possibility can be easily eliminated, and thus $b+c=9$.  Next, we write the equation from Solution 1 as $99a = 63b + 8(b+c)$, and dividing this by $9$ gives $11a = 7b+8$.  Taking both sides modulo $7$, we have $4a \equiv 1 \pmod{7}$.  Multiplying both sides by $2$ gives $a\equiv 2 \pmod{7}$, which implies $a=2$.  From here, we can find that $b=2$ and $c=7$, giving an answer of $\boxed{227}$.

 




	\item (\textit{AIME I 2022}) In isosceles trapezoid $ABCD,$ parallel bases $\overline{AB}$ and $\overline{CD}$ have lengths $500$ and $650,$ respectively, and $AD=BC=333.$ The angle bisectors of $\angle A$ and $\angle D$ meet at $P,$ and the angle bisectors of $\angle B$ and $\angle C$ meet at $Q.$ Find $PQ.$


 
 \textbf{Answer: }242



\textbf{Solution 1}

We have the following diagram:
                    \begin{figure}[H]
                        \centering
                        \includegraphics[width=0.5\textwidth]{Solution/AIME_2022_I_3_Solution_Figure_1.png}
                    \end{figure}
                    

 Let $X$ and $W$ be the points where $AP$ and $BQ$ extend to meet $CD$, and $YZ$ be the height of $\triangle AZB$. As proven in Solution 2, triangles $APD$ and $DPW$ are congruent right triangles. Therefore, $AD = DW = 333$. We can apply this logic to triangles $BCQ$ and $XCQ$ as well, giving us $BC = CX = 333$. Since $CD = 650$, $XW = DW + CX - CD = 16$.

 Additionally, we can see that $\triangle XZW$ is similar to $\triangle PQZ$ and $\triangle AZB$. We know that $\frac{XW}{AB} = \frac{16}{500}$. So, we can say that the height of the triangle $AZB$ is $500u$ while the height of the triangle $XZW$ is $16u$. After that, we can figure out the distance from $Y$ to $PQ: \frac{500+16}{2} = 258u$ and the height of triangle $PZQ: 500-258 = 242u$.

 Finally, since the ratio between the height of $PZQ$ to the height of $AZB$ is $242:500$ and $AB$ is $500$, $PQ = \boxed{242}.$

 





\textbf{Solution 2}

Extend line $PQ$ to meet $AD$ at $P'$ and $BC$ at $Q'$. The diagram looks like this:
                    \begin{figure}[H]
                        \centering
                        \includegraphics[width=0.5\textwidth]{Solution/AIME_2022_I_3_Solution_Figure_2.png}
                    \end{figure}
                    Because the trapezoid is isosceles, by symmetry $PQ$ is parallel to $AB$ and $CD$. Therefore, $\angle PAB \cong \angle APP'$ by interior angles and $\angle PAB \cong \angle PAD$ by the problem statement. Thus, $\triangle P'AP$ is isosceles with $P'P = P'A$. By symmetry, $P'DP$ is also isosceles, and thus $P'A = \frac{AD}{2}$. Similarly, the same thing is happening on the right side of the trapezoid, and thus $P'Q'$ is the midline of the trapezoid. Then, $PQ = P'Q' - (P'P + Q'Q)$.

 Since $P'P = P'A = \frac{AD}{2}, Q'Q = Q'B = \frac{BC}{2}$ and $AD = BC = 333$, we have $P'P + Q'Q = \frac{333}{2} + \frac{333}{2} = 333$. The length of the midline of a trapezoid is the average of their bases, so $P'Q' = \frac{500+650}{2} = 575$. Finally, $PQ = 575 - 333 = \boxed{242}$.

 





\textbf{Solution 3}

We have the following diagram:
                    \begin{figure}[H]
                        \centering
                        \includegraphics[width=0.5\textwidth]{Solution/AIME_2022_I_3_Solution_Figure_3.png}
                    \end{figure}
                    Extend lines $AP$ and $BQ$ to meet line $DC$ at points $W$ and $X$, respectively, and extend lines $DP$ and $CQ$ to meet $AB$ at points $Z$ and $Y$, respectively.

 Claim: quadrilaterals $AZWD$ and $BYXC$ are rhombuses.

 Proof: Since $\angle DAB + \angle ADC = 180^{\circ}$, $\angle ADP + \angle PAD = 90^{\circ}$. Therefore, triangles $APD$, $APZ$, $DPW$ and $PZW$ are all right triangles. By SAA congruence, the first three triangles are congruent; by SAS congruence, $\triangle PZW$ is congruent to the other three. Therefore, $AD = DW = WZ = AZ$, so $AZWD$ is a rhombus. By symmetry, $BYXC$ is also a rhombus.

 Extend line $PQ$ to meet $\overline{AD}$ and $\overline{BC}$ at $R$ and $S$, respectively. Because of rhombus properties, $RP = QS = \frac{333}{2}$. Also, by rhombus properties, $R$ and $S$ are the midpoints of segments $AD$ and $BC$, respectively; therefore, by trapezoid properties, $RS = \frac{AB + CD}{2} = 575$. Finally, $PQ = RS - RP - QS = \boxed{242}$.

 





\textbf{Solution 4}


                    \begin{figure}[H]
                        \centering
                        \includegraphics[width=0.5\textwidth]{Solution/AIME_2022_I_3_Solution_Figure_4.png}
                    \end{figure}
                    

 Let $X$ and $Y$ be the feet of the altitudes from $P$ and $Q$, respectively, to $AB$, and let $Z$ and $W$ be the feet of the altitudes from $P$ and $Q$, respectively, to $CD$. Side $AB$ is parallel to side $CD$, so $XYWZ$ is a rectangle with width $PQ$. Furthermore, because $CD - AB = 650-500 = 150$ and trapezoid $ABCD$ is isosceles, $WC - YB = ZD - XA = 75$. 

 Also because $ABCD$ is isosceles, $\angle ABC + \angle BCD$ is half the total sum of angles in $ABCD$, or $180^{\circ}$. Since $BQ$ and $CQ$ bisect $\angle ABC$ and $\angle BCD$, respectively, we have $\angle QBC + \angle QCB = 90^{\circ}$, so $\angle BQC = 90^{\circ}$. 

 Letting $BQ = 333k$, applying Pythagoras to $\triangle BQC$ yields $QC = 333\sqrt{1-k^2}$. We then proceed using similar triangles: $\angle BYQ = \angle BQC = 90^{\circ}$ and $\angle YBQ = \angle QBC$, so by AA similarity $YB = 333k^2$. Likewise, $\angle CWQ = \angle BQC = 90^{\circ}$ and $\angle WCQ = \angle QCB$, so by AA similarity $WC = 333(1 - k^2)$. Thus $WC + YB = 333$.

 Adding our two equations for $WC$ and $YB$ gives $2WC = 75 + 333 = 408$. Therefore, the answer is $PQ = ZW = CD - 2WC = 650 - 408 = \boxed{242}$.

 





\textbf{Solution 5}

This will be my first solution on AoPS. My apologies in advance for any errors. 

 Angle bisectors can be thought of as the locus of all points equidistant from the lines whose angle they bisect. It can thus be seen that $P$ is equidistant from  $AB, AD,$ and $CD$ and $Q$ is equidistant from  $AB, BC,$ and $CD.$ If we let the feet of the altitudes from $P$ to $AB, AD,$ and $CD$ be called $E, F,$ and $G$ respectively, we can say that $PE = PF = PG.$ Analogously, we let the feet of the altitudes from $Q$ to $AB, BC,$ and $CD$ be $H, I,$ and $J$ respectively. Thus, $QH = QI = QJ.$ Because $ABCD$ is an isosceles trapezoid, we can say that all of the altitudes are equal to each other. 

 By SA as well as SS congruence for right triangles, we find that triangles $AEP, AFP, BHQ,$ and $BIQ$ are congruent. Similarly, $DFP, DGP, CJQ,$ and $CIQ$ by the same reasoning. Additionally, $EH = GJ = PQ$ since $EHQP$ and $GJQP$ are congruent rectangles. 

 If we then let $x = AE = AF = BH = BI,$ let $y = CI = CJ = DG = DF,$ and let $z = EH = GJ = PQ,$ we can create the following system of equations with the given side length information:
 \begin{align*} 2x + z &= 500, \\ 2y + z &= 650, \\ x + y &= 333. \end{align*}
Adding the first two equations, subtracting by twice the third, and dividing by $2$ yields $z = PQ = \boxed{242}.$

 





\textbf{Solution 6}

Extend line $PQ$ to meet $AD$ at $P'$ and $BC$ at $Q'$. The diagram looks like this: 
                    \begin{figure}[H]
                        \centering
                        \includegraphics[width=0.5\textwidth]{Solution/AIME_2022_I_3_Solution_Figure_2.png}
                    \end{figure}
                    Since $\angle A + \angle D=\angle B + \angle C = 180^{\circ}$, it follows that $\angle P'AP+\angle P'DP = \angle Q'BQ + \angle Q'CQ = 90^{\circ}$. Thus, $\angle APD = \angle BQC = 90^{\circ}$, implying that $\triangle APD$ and $\triangle BQC$ are right triangles. Since $P'P$ and $Q'Q$ are medians, $P'P+Q'Q=\frac{333\times2}{2}=333$. Since $P'Q'=\frac{500+650}{2}=575$, we have $PQ+P'P+Q'Q=575$, or $PQ=575-333=\boxed{242}$.

 





\textbf{Solution 7 (Trigonometry)}

Let $PQ = x$. Note that since $AP$ bisects $\angle{A}$ and $DP$ bisects $\angle{D}$, we have 
 \[\angle{APD} = 180^{\circ}-\frac12 \angle{A}-\frac12 \angle{D}=90^{\circ}.\] Let $\angle{ADP}=\theta$. We have that $\angle{ADC} = 2\theta.$ Now, drop an altitude from $A$ to $CD$ at $E$. Notice that $DE=\frac{650-500}{2}=75$. By the definition of cosine, we have 
 \[\cos{2\theta}=1-2\cos^2{\theta}=\frac{75}{333}=\frac{25}{111} \implies \cos{\theta}=\frac{2\sqrt{1887}}{111}.\] Notice, however, that we can also apply this to $\triangle{APD}$; we have 
 \[\cos{\theta}=\frac{DP}{333} \implies DP=6\sqrt{1887}.\] By the Pythagorean Theorem, we get 
 \[AP=\sqrt{333^2-(6\sqrt{1887})^2}=3\sqrt{4773}.\] Then, drop an altitude from $P$ to $AB$ at $F$; if $AF=y$, then $PQ=x=500-2y$. Because $AP$ is an angle bisector, we see that $\angle{BAP}=\angle{DAP}=90^{\circ}-\theta$. Again, by the definition of cosine, we have 
 \[\cos{(90^{\circ}-\theta)}=\sin{\theta}=\frac{\sqrt{4773}}{111}=\frac{y}{3\sqrt{4773}} \implies y=129.\] Finally, $PQ=500-2y=\boxed{242}$.

 





\textbf{Solution 8 (Pythagoras + Similar Triangles)}


                    \begin{figure}[H]
                        \centering
                        \includegraphics[width=0.5\textwidth]{Solution/AIME_2022_I_3_Solution_Figure_4.png}
                    \end{figure}
                    

 As in solution 4, $\angle APD = 90^{\circ}$. Set $k = AX$ and $x = DP$.

 We know that $DZ = AX + \frac{DC-AB}{2}$, so $DZ = k + \frac{650-500}{2} = k + 75$.

 $\triangle DPZ \sim \triangle APD$ by AA, so we have $\frac{PD}{AD} = \frac{ZD}{PD}$, resulting in

 
 \[\frac{x}{333} = \frac{k+75}{x} \text{ (1)}\]

 $\triangle APX \sim \triangle ADP$ by AA, so we have $\frac{AP}{AD} = \frac{AX}{AP}$, resulting in

 
 \[\frac{\sqrt{333^2-x^2}}{333} = \frac{k}{\sqrt{333^2-k^2}} \text{ (2)}\]

 From $\text{(1)}$, we have $x^2 = 333k + 333(75) = 333k + 24975$. From $\text{(2)}$, we have $333^2 - x^2 = 333k$, or $x^2 = 333^2 - 333k$. Thus, $333k + 24975 = 333^2 - 333k$. Solving for $k$ yields $k = 129$.

 By symmetry, $YB = AX = 129$. Thus, $PQ = XY = AB - 2AX = 500 - 2(129) = \boxed{242}$.

 





\textbf{Solution 9 (Really Quick)}

Extend $PQ$. Let the endpoints of this extension be $X$ and $Y$, ($X$ is closer to $P$ and $Y$ is closer to $Q$). Now let $PX = QY = a$. So, $PQ = 575-2a$. Now, by transversals, $\angle{PAB} = \angle{XPA}$. So, $\triangle{XAP}$ is isosceles with $XA = XP = a$. Now, lets look at $XD$. $XD = 333-a$. Also, by transversals, $\angle{DPX} = \angle{PDC}$. Thus, we have $XD = XP$. Aha! So, $333-a = a \Longrightarrow a = \dfrac{333}{2}$. Plugging this into our equation for $PQ$, we get $575-333 = \boxed{242}$

 





\textbf{Solution 10 (Fast Trigonometric Solve)}

Noticing immediately that $\Delta APD$ and $\Delta BQC$ are right triangles due to the angle bisector conditions, we are motivated to use trigonometry and half-angles. An obvious first step is to compute the cosine of $\angle{D}$, which we can drop the altitude from $A$ and find $\cos(\angle ADC) = \frac{25}{111}$. We can easily then see that by half angles, $\cos(\angle ADP) = \cos(\angle PDC) = \sqrt{\frac{68}{111}}$. 

 Motivated by right triangle trigonometry, we drop the altitude from $P$ to $CD$ to form point $X$. By our derived cosines, $\frac{DP}{AD} = \frac{DE}{PD} = \sqrt{\frac{68}{111}}$, so we can do some mathematics to find $DE = AD \cdot \frac{68}{111} = \frac{333 \cdot 68}{111} = 204$. To find $PQ$, we can drop the altitude from $Q$ similarly to the construction of $X$ to form point $Y$, and $PQ = XY = CD - DX - CY = CD = 2DX = 650-2\cdot204 = \boxed{242}$

 





\textbf{Solution 11 (Continuing on)}

We notice that it is a right triangle. As mentioned by solution 2 above, we can use parallel lines to get that the length PP` is just the median of the triangle and similar for QQ` so these are just 333/2 which sum up to 333 so we have 575 - 333 = 242.

 





\textbf{Solution 12 (basic trig)}

Notice that the desired length is equal to $500-2\overline{AP}\cos (\angle PAB)$ as well as $650-2\overline{DP}\cos (\angle PDC)$. Obviously $\angle APD = 90^\circ$, so $\cos(\angle PAB)=\frac{\overline AP}{333}$ and $\cos(\angle PDC)=\frac{\overline DP}{333}$.

 Hence $500-\frac{2\overline{AP}^2}{333}=650-\frac{2\overline{DP}^2}{333}$, and $\overline{DP}^2-\overline{AP^2}=333\cdot 75$. But obviously $\overline{DP}^2+\overline{AP^2}=333\cdot 333$, hence $2\overline{DP}^2=408*333$. Then it is clear that $650-\frac{2\overline{DP}^2}{333}=650-408=\boxed{242}$.

 -benjamintontungtungtungsahur

 




	\item (\textit{AIME I 2022}) Let $w = \dfrac{\sqrt{3} + i}{2}$ and $z = \dfrac{-1 + i\sqrt{3}}{2},$ where $i = \sqrt{-1}.$ Find the number of ordered pairs $(r,s)$ of positive integers not exceeding $100$ that satisfy the equation $i \cdot w^r = z^s.$


 
 \textbf{Answer: }834



\textbf{Solution 1}

We rewrite $w$ and $z$ in polar form:
 \begin{align*} w &= e^{i\cdot\frac{\pi}{6}}, \\ z &= e^{i\cdot\frac{2\pi}{3}}. \end{align*}
The equation $i \cdot w^r = z^s$ becomes
 \begin{align*} e^{i\cdot\frac{\pi}{2}} \cdot \left(e^{i\cdot\frac{\pi}{6}}\right)^r &= \left(e^{i\cdot\frac{2\pi}{3}}\right)^s \\ e^{i\left(\frac{\pi}{2}+\frac{\pi}{6}r\right)} &= e^{i\left(\frac{2\pi}{3}s\right)} \\ \frac{\pi}{2}+\frac{\pi}{6}r &= \frac{2\pi}{3}s+2\pi k \\ 3+r &= 4s+12k \\ 3+r &= 4(s+3k). \end{align*}
for some integer $k.$ 

 Since $4\leq 3+r\leq 103$ and $4\mid 3+r,$ we conclude that
 \begin{align*} 3+r &\in \{4,8,12,\ldots,100\}, \\ s+3k &\in \{1,2,3,\ldots,25\}. \end{align*}
Note that the values for $s+3k$ and the values for $r$ have one-to-one correspondence.

 We apply casework to the values for $s+3k:$

 
            \begin{enumerate}
                 \item $s+3k\equiv0\pmod{3}$There are $8$ values for $s+3k,$ so there are $8$ values for $r.$ It follows that $s\equiv0\pmod{3},$ so there are $33$ values for $s.$ 

 There are $8\cdot33=264$ ordered pairs $(r,s)$ in this case.  

 \item $s+3k\equiv1\pmod{3}$There are $9$ values for $s+3k,$ so there are $9$ values for $r.$ It follows that $s\equiv1\pmod{3},$ so there are $34$ values for $s.$ 

 There are $9\cdot34=306$ ordered pairs $(r,s)$ in this case.  

 \item $s+3k\equiv2\pmod{3}$There are $8$ values for $s+3k,$ so there are $8$ values for $r.$ It follows that $s\equiv2\pmod{3},$ so there are $33$ values for $s.$ 

 There are $8\cdot33=264$ ordered pairs $(r,s)$ in this case.

 
            \end{enumerate}
            Together, the answer is $264+306+264=\boxed{834}.$

 





\textbf{Solution 2}

First we recognize that $w = \operatorname{cis}(30^{\circ})$ and $z = \operatorname{cis}(120^{\circ})$ because the cosine and sine sums of those angles give the values of $w$ and $z$, respectively. By De Moivre's theorem, $\operatorname{cis}(\theta)^n = \operatorname{cis}(n\theta)$. When you multiply by $i$, we can think of that as rotating the complex number $90^{\circ}$ counterclockwise in the complex plane. Therefore, by the equation we know that $30r + 90$ and $120s$ land on the same angle.

 This means that
 \[30r + 90 \equiv 120s \pmod{360},\]which we can simplify to 
 \[r+3 \equiv 4s \pmod{12}.\]Notice that this means that $r$ cycles by $12$ for every value of $s$. This is because once $r$ hits $12$, we get an angle of $360^{\circ}$ and the angle laps onto itself again. By a similar reasoning, $s$ laps itself every $3$ times, which is much easier to count. By listing the possible values out, we get the pairs $(r,s)$:
 \[\begin{array}{cccccccc} (1,1) & (5,2) & (9,3) & (13,1) & (17,2) & (21,3) & \ldots & (97,1) \\ (1,4) & (5,5) & (9,6) & (13,4) & (17,5) & (21,6) & \ldots & (97,4) \\  (1,7) & (5,8) & (9,9) & (13,7) & (17,8) & (21,9) & \ldots & (97,7) \\ [-1ex] \vdots & \vdots & \vdots & \vdots & \vdots & \vdots & \vdots & \vdots \\ (1,100) & (5,98) & (9,99) & (13,100) & (17,98) & (21,99) & \ldots & (97,100) \end{array}\]We have $25$ columns in total: $34$ values for the first column, $33$ for the second, $33$ for the third, and then $34$ for the fourth, $33$ for the fifth, $33$ for the sixth, etc. Therefore, this cycle repeats every $3$ columns and our total sum is $(34+33+33) \cdot 8 + 34 = 100 \cdot 8 + 34 = \boxed{834}$.

 \href{https://mathproblemsolvingskills.wordpress.com/}{MathProblemSolvingSkills.com}

 




	\item (\textit{AIME I 2022}) A straight river that is $264$ meters wide flows from west to east at a rate of $14$ meters per minute. Melanie and Sherry sit on the south bank of the river with Melanie a distance of $D$ meters downstream from Sherry. Relative to the water, Melanie swims at $80$ meters per minute, and Sherry swims at $60$ meters per minute. At the same time, Melanie and Sherry begin swimming in straight lines to a point on the north bank of the river that is equidistant from their starting positions. The two women arrive at this point simultaneously. Find $D.$


 
 \textbf{Answer: }550



\textbf{Solution 1 (Euclidean)}

Define $m$ as the number of minutes they swim for.

 Let their meeting point be $A$. Melanie is swimming against the current, so she must aim upstream from point $A$, to compensate for this; in particular, since she is swimming for $m$ minutes, the current will push her $14m$ meters downstream in that time, so she must aim for a point $B$ that is $14m$ meters upstream from point $A$. Similarly, Sherry is swimming downstream for $m$ minutes, so she must also aim at point $B$ to compensate for the flow of the current.

 If Melanie and Sherry were to both aim at point $B$ in a \textit{currentless} river with the same dimensions, they would still both meet at that point simultaneously. Since there is no current in this scenario, the distances that Melanie and Sherry travel, respectively, are $80m$ and $60m$ meters. We can draw out this new scenario, with the dimensions that we have:
                    \begin{figure}[H]
                        \centering
                        \includegraphics[width=0.5\textwidth]{Solution/AIME_2022_I_5_Solution_Figure_1.png}
                    \end{figure}
                    (While it is indeed true that the triangle above with side lengths $60m$, $80m$ and $D$ is a right triangle, we do not know this yet, so we cannot assume this based on the diagram.)

 By the Pythagorean Theorem, we have
 \begin{align*} 264^{2} + \left( \frac{D}{2} - 14m \right) ^{2} &= 3600m^{2} \\ 264^{2} + \left( \frac{D}{2} + 14m \right) ^{2} &= 6400m^{2}. \end{align*}


 Subtracting the first equation from the second gives us $28Dm = 2800m^{2}$, so $D = 100m$. Substituting this into our first equation, we have that
 \begin{align*}264^{2} + 36^{2} m^{2} &= 60^{2}m^{2} \\ 264^{2} &= 96 \cdot 24 \cdot m^{2} \\ 11^{2} &= 4 \cdot m^{2} \\ m &= \frac{11}{2}. \end{align*}


 So $D = 100m = \boxed{550}$.

 





\textbf{Solution 2 (Euclidean)}

<i>\textbf{Claim}</i>

 
                    \begin{figure}[H]
                        \centering
                        \includegraphics[width=0.5\textwidth]{Solution/AIME_2022_I_5_Solution_Figure_2.png}
                    \end{figure}
                    Median $AM$ and altitude $AH$ are drawn in triangle $ABC$.$AB = c, AC = b < c, BC = a$ are known. Let's denote $MH = x$.

 Prove that
 \begin{align*}2ax = c^{2} - b^{2}\end{align*}


 <i>\textbf{Proof}</i>
 \[BH + CH = a,\]
 \begin{align*} BH^{2} - CH^{2} = c^{2} - b^{2}\implies BH - CH &= \frac{c^{2} - b^{2}} {a},\end{align*}

 \[BH =  \frac{c^{2} - b^{2}}{2a} +  \frac{a}{2},\]
 \begin{align*}MH = BH - BM  &= \frac{c^{2} - b^{2}} {2a}.\end{align*}


 <i>\textbf{Solution}</i>

 In the coordinate system associated with water, the movement is described by the scheme in the form of a triangle, the side on which Melanie floats is $80t$, where t is the time of  Melanie's movement, the side along which Sherry floats is $60t$.

 The meeting point floated away at a distance of $14t$ from the midpoint between the starting points of Melanie and Sherry.

 In the notation of the <i>\textbf{Claim}</i>,
 \begin{align*}  c = 80t, b = 60t, x = 14t \implies a = \frac{(80t)^2-(60t)^2}{2 \cdot 14t}=\frac{20^2}{4}\cdot \frac{16-9}{7}t = 100t.\end{align*}
Hence,
 \begin{align*}  AH = \sqrt{BC^2-BH^2}= \sqrt{(80t)^2-(50t+14t)^2}=16t \cdot \sqrt{5^2-4^2}= 48t = 264 \implies t = 5.5.\end{align*}

 \[D = a = 100t = \boxed{550}\]\textbf{vladimir.shelomovskii@gmail.com, vvsss}

 





\textbf{Solution 3 (Vectors)}

We have the following diagram:
                    \begin{figure}[H]
                        \centering
                        \includegraphics[width=0.5\textwidth]{Solution/AIME_2022_I_5_Solution_Figure_3.png}
                    \end{figure}
                    Since Melanie and Sherry swim for the same distance and the same amount of time, they swim at the same net speed.

 Let $x$ and $y$ be some positive numbers. We have the following table:
 \[\begin{array}{c||c|c|c}     & \textbf{Net Velocity Vector (m/min)} & \textbf{Natural Velocity Vector (m/min)} & \textbf{Natural Speed (m/min)} \\ \hline \hline &&& \\ [-2.25ex] \textbf{Melanie} & \langle -x,y\rangle & \langle -x-14,y\rangle & 80  \\ \hline   &&& \\ [-2.25ex] \textbf{Sherry} & \langle x,y\rangle & \langle x-14,y\rangle & 60 \end{array}\]Recall that $|\text{velocity}|=\text{speed},$ so
 \begin{align*} (-x-14)^2 + y^2 &= 80^2, &&(1) \\ (x-14)^2 + y^2 &= 60^2. &&(2) \end{align*}
We subtract $(2)$ from $(1)$ to get $56x=2800,$ from which $x=50.$ Substituting this into either equation, we have $y=48.$

 It follows that Melanie and Sherry both swim for $264\div y=5.5$ minutes. Therefore, the answer is 
 \[D=2x\cdot5.5=\boxed{550}.\]~MRENTHUSIASM

 





\textbf{Solution 4 (Vectors)}

We can break down movement into two components: the $x$-component and the $y$-component. Suppose that Melanie travels a distance of $a$ in the $x$-direction and a distance of $c$ in the $y$-direction in one minute when there is no current. Similarly, suppose that Sherry travels a distance of $a$ in the $x$-direction but a distance of $b$ in the $y$-direction in one minute when there is no current. The current only affects the $x$-components because it goes in the $x$-direction. 

 
                    \begin{figure}[H]
                        \centering
                        \includegraphics[width=0.5\textwidth]{Solution/AIME_2022_I_5_Solution_Figure_4.png}
                    \end{figure}
                    

 Now, note that $a^2 + b^2 = 60^2$ because Sherry travels 60 meters in a minute. Thus, $a^2 + c^2 = 80^2$ because Melanie travels 80 meters in a minute. Also, the distance they travel with the current must be the same in one minute because they reach the point equidistant from them at the same time. That means $b + 14 = c - 14$ or $b = c - 28$. So now we can plug that into the two equations to get:
 \begin{align*} a^2 + c^2 &= 80^2, \\ a^2 + (c-28)^2 &= 60^2. \end{align*}
We can solve the system of equations to get $a = 48$ and $c = 64$. From this, we can figure out that it must've taken them $5.5$ minutes to get to the other side, because $264/48 = 5.5$. This means that there are $5.5$ lengths of $48$ in each person's travel. Also, $D$ must be equal to $11(b+14) = 11(c-14)$ because there are $(5.5)2 = 11$ lengths of $b-14$ between them, $5.5$ on each person's side. Since $c = 64$, we have $c-14 = 50$, so the answer is 
 \[D=11\cdot50=\boxed{550}.\]~Curious\_crow

 




	\item (\textit{AIME I 2022}) Find the number of ordered pairs of integers $(a,b)$ such that the sequence 
 \[3,4,5,a,b,30,40,50\] is strictly increasing and no set of four (not necessarily consecutive) terms forms an arithmetic progression.


 
 \textbf{Answer: }228



\textbf{Solution 1}

Since $3,4,5,a$ and $3,4,5,b$ cannot be an arithmetic progression, $a$ or $b$ can never be $6$. Since $b, 30, 40, 50$ and $a, 30, 40, 50$ cannot be an arithmetic progression, $a$ and $b$ can never be $20$. Since $a < b$, there are ${24 - 2 \choose 2} = 231$ ways to choose $a$ and $b$ with these two restrictions in mind.

 However, there are still specific invalid cases counted in these $231$ pairs $(a,b)$. Since
 \[3,5,a,b\]cannot form an arithmetic progression, $\underline{(a,b) \neq (7,9)}$.
 \[a,b,30,50\]cannot be an arithmetic progression, so $(a,b) \neq (-10,10)$; however, since this pair was not counted in our $231$, we do not need to subtract it off.
 \[3,a,b,30\]cannot form an arithmetic progression, so $\underline{(a,b) \neq (12,21)}$.
 \[4, a, b, 40\]cannot form an arithmetic progression, so $\underline{(a,b) \neq (16,28)}$.
 \[5, a,b, 50\]cannot form an arithmetic progression, $(a,b) \neq 20, 35$; however, since this pair was not counted in our $231$ (since we disallowed $a$ or $b$ to be $20$), we do not to subtract it off.

 Also, the sequences $(3,a,b,40)$, $(3,a,b,50)$, $(4,a,b,30)$, $(4,a,b,50)$, $(5,a,b,30)$ and $(5,a,b,40)$ will never be arithmetic, since that would require $a$ and $b$ to be non-integers.

 So, we need to subtract off $3$ progressions from the $231$ we counted, to get our final answer of $\boxed{228}$.

 





\textbf{Solution 2 (Rigorous)}

We will follow the solution from earlier in a rigorous manner to show that there are no other cases missing.

 We recognize that an illegal sequence (defined as one that we subtract from our 231) can never have the numbers {3, 4} and {4,5} because we have not included a 6 in our count. Similarly, sequences with {30,40} and {40,50} will not give us any subtractions because those sequences must all include a 20. Let's stick with the lower ones for a minute: if we take them two at a time, then {3,5} will give us the subtraction of 1 sequence {3,5,7,9}. We have exhausted all pairs of numbers we can take, and if we take the triplet of single digit numbers, the only possible sequence must have a 6, which we already don't count. Therefore, we subtract $\textbf{1}$ from the count of illegal sequences with any of the single-digit numbers and none of the numbers 30,40,50.(Note if we take only 1 at a time, there will have to be 3 of $a, b$, which is impossible.)

 If we have the sequence including {30,50}, we end up having negative values, so these do not give us any subtractions, and the triplet {30,40,50} gives us a 20. Hence by the same reasoning as earlier, we have 0 subtractions from the sequences with these numbers and none of the single digit numbers {3,4,5}.

 Finally, we count the sequences that are something like (one of 3,4,5,), $a, b$, (one of 30, 40, 50). If this is to be the case, then let $a$ be the starting value in the sequence. The sequence will be $a, a+d, a+2d, a+3d$; We see that if we subtract the largest term by the smallest term we have $3d$, so the subtraction of one of (30,40,50) and one of (3,4,5) must be divisible by 3. Therefore the only sequences possible are $3,a,b,30; 4,a,b,40; 5,a,b,50$. Of these, only the last is invalid because it gives $b = 35$, larger than our bounds $6<a<b<30$. Therefore, we subtract $\textbf{2}$ from this case.

 Our final answer is $231 - 1 - 2 = \boxed{228}$

 





\textbf{Solution 3}

Denote $S = \left\{ (a, b) : 6 \leq a < b \leq 29 \right\}$.

 Denote by $A$ a subset of $S$, such that there exists an arithmetic sequence that has 4 terms and includes $a$ but not $b$.

 Denote by $B$ a subset of $S$, such that there exists an arithmetic sequence that has 4 terms and includes $b$ but not $a$.

 Hence, $C$ is a subset of $S$, such that there exists an arithmetic sequence that has 4 terms and includes both $a$ and $b$.

 Hence, this problem asks us to compute
 \[ | S | - \left( | A | + | B | + | C | \right) . \]

 First, we compute $| S |$.

 We have $| S | = \binom{29 - 6 + 1}{2} = \binom{24}{2} = 276$.

 Second, we compute $| A |$.

 $\textbf{Case 1}$: $a = 6$.

 We have $b = 8 , \cdots , 19, 21, 22, \cdots, 29$.Thus, the number of solutions is 22.

 $\textbf{Case 2}$: $a = 20$.

 We have $b = 21, 22, \cdots , 29$.Thus, the number of solutions is 9.

 Thus, $| A | = 22 + 9 = 31$.

 Third, we compute $| B |$.

 In $B$, we have $b = 6, 20$. However, because $6 \leq a < b$, we have $b \geq 7$.Thus, $b = 20$.

 This implies $a = 7, 8, 9, 11, 12, \cdots , 19$. Note that $(a, b)=(10, 20)$ belongs in $C$.

 Thus, $| B | = 12$.

 Fourth, we compute $| C |$.

 $\textbf{Case 1}$: In the arithmetic sequence, the two numbers beyond $a$ and $b$ are on the same side of $a$ and $b$.

 Hence, $(a, b) = (6 , 7), (7, 9) , (10, 20)$.Therefore, the number solutions in this case is 3.

 $\textbf{Case 2}$: In the arithmetic sequence, the two numbers beyond $a$ and $b$ are on the opposite sides of $a$ and $b$.

 $\textbf{Case 2.1}$: The arithmetic sequence is $3, a, b, 30$.

 Hence, $(a, b) = (12, 21)$.

 $\textbf{Case 2.2}$: The arithmetic sequence is $4, a, b, 40$.

 Hence, $(a, b) = (16, 28)$.

 $\textbf{Case 2.3}$: The arithmetic sequence is $5, a, b, 50$.

 Hence, $(a, b) = (20, 35)$. However, the sequence $... 20, 35, 30, 40, 50$ is not strictly increasing.

 Putting two cases together, $| C | = 65.$

 Therefore,
 \[| S | - \left( | A | + | B | + | C | \right) = 276 - \left( 31 + 12 + 5 \right) = \boxed{228}.\]

 





\textbf{Solution 4}

divide cases into $7\leq a<20; 21\leq a\leq28$.(Notice that $a$ can't be equal to $6,20$, that's why I divide them into two parts.There are three cases that arithmetic sequence forms: $3,12,21,30;4,16,28,40;3,5,7,9$.(NOTICE that $5,20,35,50$ IS NOT A VALID SEQUENCE!)So when $7\leq a<20$, there are $10+11+12+...+22-3-13=192$ possible ways( 3 means the arithmetic sequence and 13 means there are 13 "a" s and b cannot be 20)

 When $21\leq a \leq 28$, there are $1+2+\cdots+8=36$ ways.

 In all, there are $192+36=\boxed{228}$ possible sequences.

 




	\item (\textit{AIME I 2022}) Let $a,b,c,d,e,f,g,h,i$ be distinct integers from $1$ to $9.$ The minimum possible positive value of 
 \[\dfrac{a \cdot b \cdot c - d \cdot e \cdot f}{g \cdot h \cdot i}\] can be written as $\frac{m}{n},$ where $m$ and $n$ are relatively prime positive integers. Find $m+n.$


 
 \textbf{Answer: }289



\textbf{Solution 1 (Optimization)}

To minimize a positive fraction, we minimize its numerator and maximize its denominator. It is clear that $\frac{a \cdot b \cdot c - d \cdot e \cdot f}{g \cdot h \cdot i} \geq \frac{1}{7\cdot8\cdot9}.$

 If we minimize the numerator, then $a \cdot b \cdot c - d \cdot e \cdot f = 1.$ Note that $a \cdot b \cdot c \cdot d \cdot e \cdot f = (a \cdot b \cdot c) \cdot (a \cdot b \cdot c - 1) \geq 6! = 720,$ so $a \cdot b \cdot c \geq 28.$ It follows that $a \cdot b \cdot c$ and $d \cdot e \cdot f$ are consecutive composites with prime factors no other than $2,3,5,$ and $7.$ The smallest values for $a \cdot b \cdot c$ and $d \cdot e \cdot f$ are $36$ and $35,$ respectively. So, we have $\{a,b,c\} = \{2,3,6\}, \{d,e,f\} = \{1,5,7\},$ and $\{g,h,i\} = \{4,8,9\},$ from which $\frac{a \cdot b \cdot c - d \cdot e \cdot f}{g \cdot h \cdot i} = \frac{1}{288}.$

 If we do not minimize the numerator, then $a \cdot b \cdot c - d \cdot e \cdot f > 1.$ Note that $\frac{a \cdot b \cdot c - d \cdot e \cdot f}{g \cdot h \cdot i} \geq \frac{2}{7\cdot8\cdot9} > \frac{1}{288}.$

 Together, we conclude that the minimum possible positive value of $\frac{a \cdot b \cdot c - d \cdot e \cdot f}{g \cdot h \cdot i}$ is $\frac{1}{288}.$ Therefore, the answer is $1+288=\boxed{289}.$

 





\textbf{Solution 2 (Bash)}

Obviously, to find the correct answer, we need to get the largest denominator with the smallest numerator. 

 To bash efficiently, we can start out with $7\cdot8\cdot9$ as our denominator. This, however, leaves us with the numbers $1, 2, 3, 4, 5,$ and $6$ left. The smallest we can make out of this is $1\cdot5\cdot6 - 2\cdot3\cdot4 = 30 - 24 = 6$. When simplified, it gives us $\frac{1}{84}$, which gives a small answer of $85$. Obviously there are larger answers than this.

 After the first bash, we learn to bash even more efficiently, we can consider <u>both</u> the numerator and the denominator when guessing. We know the numerator has to be extremely small while still having a large denominator. When bashing, we soon find out the couple $(a,b,c)=(2,3,6)$ and $(d,e,f)=(1,5,7)$.

 This gives us a numerator of $36-35=1$, which is by far the smallest yet. With the remaining numbers $4, 8,$ and $9$, we get $\frac{36-35}{4\cdot8\cdot9}=\frac{1}{288}$.

 Finally, we add up our numerator and denominator: The answer is $1+288=\boxed{289}$.

 





\textbf{Solution 3 (Educated Trial and Error)}

To minimize the numerator, we must have $abc - def = 1$. Thus, one of these products must be odd and the other must be even. The odd product must consist of only odd numbers. The smallest such value $(d, e, f) = (1, 3, 5)$ cannot result in a difference of $1$, and the next smallest product, $(d, e, f) = (1, 3, 7)$ cannot either, but $(d, e, f) = (1, 5, 7)$ can if $(a, b, c) = (2, 3, 6)$. Thus, the denominator must be $(g, h, i) = (4, 8, 9)$, and the smallest fraction possible is $\dfrac{36 - 35}{288} = \dfrac{1}{288}$, making the answer $1 + 288 = \boxed{289}$.

 




	\item (\textit{AIME I 2022}) Equilateral triangle $\triangle ABC$ is inscribed in circle $\omega$ with radius $18.$ Circle $\omega_A$ is tangent to sides $\overline{AB}$ and $\overline{AC}$ and is internally tangent to $\omega.$ Circles $\omega_B$ and $\omega_C$ are defined analogously. Circles $\omega_A,$ $\omega_B,$ and $\omega_C$ meet in six points---two points for each pair of circles. The three intersection points closest to the vertices of $\triangle ABC$ are the vertices of a large equilateral triangle in the interior of $\triangle ABC,$ and the other three intersection points are the vertices of a smaller equilateral triangle in the interior of $\triangle ABC.$ The side length of the smaller equilateral triangle can be written as $\sqrt{a} - \sqrt{b},$ where $a$ and $b$ are positive integers. Find $a+b.$


 
 \textbf{Answer: }378



\textbf{Solution 1 (Coordinate Geometry)}

We can extend $AB$ and $AC$ to $B'$ and $C'$ respectively such that circle $\omega_A$ is the incircle of $\triangle AB'C'$.
                    \begin{figure}[H]
                        \centering
                        \includegraphics[width=0.5\textwidth]{Solution/AIME_2022_I_8_Solution_Figure_1.png}
                    \end{figure}
                    Since the diameter of the circle is the height of this triangle, the height of this triangle is $36$. We can use inradius or equilateral triangle properties to get the inradius of this triangle is $12$ (The incenter is also a centroid in an equilateral triangle, and the distance from a side to the centroid is a third of the height). Therefore, the radius of each of the smaller circles is $12$.

 Let $O=\omega$ be the center of the largest circle. We will set up a coordinate system with $O$ as the origin. The center of $\omega_A$ will be at $(0,-6)$ because it is directly beneath $O$ and is the length of the larger radius minus the smaller radius, or $18-12 = 6$. By rotating this point $120^{\circ}$ around $O$, we get the center of $\omega_B$. This means that the magnitude of vector $\overrightarrow{O\omega_B}$ is $6$ and is at a $30$ degree angle from the horizontal. Therefore, the coordinates of this point are $(3\sqrt{3},3)$ and by symmetry the coordinates of the center of $\omega_C$ is $(-3\sqrt{3},3)$.

 The upper left and right circles intersect at two points, the lower of which is $X$. The equations of these two circles are:
 \begin{align*} (x+3\sqrt3)^2 + (y-3)^2 &= 12^2, \\ (x-3\sqrt3)^2 + (y-3)^2 &= 12^2. \end{align*}
We solve this system by subtracting to get $x = 0$. Plugging back in to the first equation, we have $(3\sqrt{3})^2 + (y-3)^2 = 144 \implies (y-3)^2 = 117 \implies y-3 = \pm \sqrt{117} \implies y = 3 \pm \sqrt{117}$. Since we know $X$ is the lower solution, we take the negative value to get $X = (0,3-\sqrt{117})$.

 We can solve the problem two ways from here. We can find $Y$ by rotation and use the distance formula to find the length, or we can be somewhat more clever. We notice that it is easier to find $OX$ as they lie on the same vertical, $\angle XOY$ is $120$ degrees so we can make use of $30-60-90$ triangles, and $OX = OY$ because $O$ is the center of triangle $XYZ$. We can draw the diagram as such:
                    \begin{figure}[H]
                        \centering
                        \includegraphics[width=0.5\textwidth]{Solution/AIME_2022_I_8_Solution_Figure_2.png}
                    \end{figure}
                    Note that $OX = OY = \sqrt{117} - 3$. It follows that
 \begin{align*} XY &= 2 \cdot \frac{OX\cdot\sqrt{3}}{2} \\ &= OX \cdot \sqrt{3} \\ &= (\sqrt{117}-3) \cdot \sqrt{3} \\ &= \sqrt{351}-\sqrt{27}. \end{align*}
Finally, the answer is $351+27 = \boxed{378}$.

 





\textbf{Solution 2 (Euclidean Geometry)}


                    \begin{figure}[H]
                        \centering
                        \includegraphics[width=0.5\textwidth]{Solution/AIME_2022_I_8_Solution_Figure_3.png}
                    \end{figure}
                    For equilateral triangle with side length $l$, height $h$, and circumradius $r$, there are relationships: $h = \frac{\sqrt{3}}{2} l$, $r = \frac{2}{3} h = \frac{\sqrt{3}}{3} l$, and $l = \sqrt{3}r$.

 There is a lot of symmetry in the figure. The radius of the big circle $\odot \omega$ is $R = 18$, let the radius of the small circles $\odot \omega_A$, $\odot \omega_B$, $\odot \omega_C$ be $r$.

 We are going to solve this problem in $3$ steps:

 $\textbf{Step 1:}$ 

 We have $\triangle A \omega_A D$ is a $30-60-90$ triangle, and $A \omega_A = 2 \cdot \omega_A D$, $A \omega_A = 2R-r$ ($\odot \omega$ and $\odot \omega_A$ are tangent), and $\omega_A D = r$. So, we get $2R-r = 2r$ and $r = \frac{2}{3} \cdot R = 12$.

 Since $\odot \omega$ and $\odot \omega_A$ are tangent, we get $\omega \omega_A = R - r = \frac{1}{3} \cdot R = 6$.

 Note that $\triangle \omega_A \omega_B \omega_C$ is an equilateral triangle, and $\omega$ is its center, so $\omega_B \omega_C = \sqrt{3} \cdot \omega \omega_A = 6 \sqrt{3}$.

 $\textbf{Step 2:}$ 

 Note that $\triangle \omega_C E X$ is an isosceles triangle, so 
 \[EX = 2 \sqrt{(\omega_C E)^2 - \left(\frac{\omega_B \omega_C}{2}\right)^2} = 2 \sqrt{r^2 - \left(\frac{\omega_B \omega_C}{2}\right)^2} = 2 \sqrt{12^2 - (3 \sqrt{3})^2} = 2 \sqrt{117}.\]

 $\textbf{Step 3:}$ 

 In $\odot \omega_C$, Power of a Point gives $\omega X \cdot \omega E = r^2 - (\omega_C \omega)^2$ and $\omega E = EX - \omega X = 2\sqrt{117} - \omega X$.

 It follows that $\omega X \cdot (2\sqrt{117} - \omega X) = 12^2 - 6^2$. We solve this quadratic equation: $\omega X = \sqrt{117} - 3$.

 Since $\omega X$ is the circumradius of equilateral $\triangle XYZ$, we have $XY = \sqrt{3} \cdot \omega X = \sqrt{3} \cdot (\sqrt{117} - 3) = \sqrt{351}-\sqrt{27}$.

 Therefore, the answer is $351+27 = \boxed{378}$.

 





\textbf{Solution 3 (Simple Geometry)}


                    \begin{figure}[H]
                        \centering
                        \includegraphics[width=0.5\textwidth]{Solution/AIME_2022_I_8_Solution_Figure_4.png}
                    \end{figure}
                    Let $O$ be the center, $R = 18$ be the radius, and $CC'$ be the diameter of $\omega.$ Let $r$  be the radius, $E,D,F$ are the centers of $\omega_A, \omega_B,\omega_C.$ Let $KGH$ be the desired triangle with side  $x.$We find $r$ using 
 \[CC' = 2R = C'K + KC = r + \frac{r}{\sin 30^\circ} = 3r.\]
 \[r = \frac{2R}{3} = 12.\]
 \[OE = R – r = 6.\]Triangles $\triangle DEF$ and $\triangle KGH$ – are equilateral triangles with a common center $O,$ therefore in the triangle  $OEH$ $OE = 6, \angle EOH = 120^\circ, OH = \frac{x}{\sqrt3}.$

 We apply the Law of Cosines to $\triangle OEH$ and get
 \[OE^2 + OH^2 + OE \cdot OH = EH^2.\]
 \[6^2 + \frac{x^2}{3} + \frac{6x}{\sqrt3} = 12^2.\]
 \[x^2 + 6x \sqrt{3} = 324\]
 \[x= \sqrt{351} - \sqrt{27}  \implies 351 + 27 = \boxed {378}\]

 \textbf{vladimir.shelomovskii@gmail.com, vvsss}

 





\textbf{Solution 4 (Mixtilinear Incircles)}

Let $O$ be the center of $\omega$, $X$ be the intersection of $\omega_B,\omega_C$ further from $A$, and $O_A$ be the center of $\omega_A$. Define $Y, Z, O_B, O_C$ similarly. It is well-known that the $A$-mixtilinear inradius $R_A$ is $\frac{r}{\cos^2\left(\frac{\angle A}{2}\right)} = \frac{9}{\cos^2\left(30^{\circ}\right)} = 12$, so in particular this means that $OO_B = 18 -  R_B = 6 = OO_C$. Since $\angle O_BOO_C = \angle BOC = 120^\circ$, it follows by Law of Cosines on $\triangle OO_BO_C$ that $O_BO_C = 6\sqrt{3}$. Then the Pythagorean theorem gives that the altitude of $O_BO_CX$ is $\sqrt{117}$, so $OY = OX = \text{dist}(X, YZ) - \text{dist}(O, YZ) = \sqrt{117} - 3$ and $YZ = \frac{O_BO_C\cdot OY}{OO_B} = \frac{6\sqrt{3}(\sqrt{117} - 3)}{6}=\sqrt{351} - \sqrt{27}$ so the answer is $351 + 27 = \boxed{378}$.

 




	\item (\textit{AIME I 2022}) Ellina has twelve blocks, two each of red ($\textbf{R}$), blue ($\textbf{B}$), yellow ($\textbf{Y}$), green ($\textbf{G}$), orange ($\textbf{O}$), and purple ($\textbf{P}$). Call an arrangement of blocks $\textit{even}$ if there is an even number of blocks between each pair of blocks of the same color. For example, the arrangement

 \[\textbf{R B B Y G G Y R O P P O}\]
is even. Ellina arranges her blocks in a row in random order. The probability that her arrangement is even is $\frac{m}{n},$ where $m$ and $n$ are relatively prime positive integers. Find $m+n.$


 
 \textbf{Answer: }247



\textbf{Solution 1}

Consider this position chart: 
 \[\textbf{1 2 3 4 5 6 7 8 9 10 11 12}\]Since there has to be an even number of spaces between each pair of the same color, spots $1$, $3$, $5$, $7$, $9$, and $11$ contain some permutation of all $6$ colored balls. Likewise, so do the even spots, so the number of even configurations is $6! \cdot 6!$ (after putting every pair of colored balls in opposite parity positions, the configuration can be shown to be even). This is out of $\frac{12!}{(2!)^6}$ possible arrangements, so the probability is: 
 \[\frac{6!\cdot6!}{\frac{12!}{(2!)^6}} = \frac{6!\cdot2^6}{7\cdot8\cdot9\cdot10\cdot11\cdot12} = \frac{2^4}{7\cdot11\cdot3} = \frac{16}{231},\]which is in simplest form. So, $m + n = 16 + 231 = \boxed{247}$.

 





\textbf{Solution 2}

We can simply use constructive counting. First, let us place the red blocks; choose the first slot in $12$ ways, and the second in $6$ ways, because the number is cut in half due to the condition in the problem. This gives $12 \cdot 6$ ways to place the red blocks. Similarly, there are $10 \cdot 5$ ways to place the blue blocks, and so on, until there are $2 \cdot 1$ ways to place the purple blocks. Thus, the probability is 
 \[\frac{12 \cdot 6 \cdot 10 \cdot 5 \cdot 8 \cdot 4 \cdot 6 \cdot 3 \cdot 4 \cdot 2 \cdot 2 \cdot 1}{12!}=\frac{16}{231},\] and the desired answer is $16+231=\boxed{247}$.

 





\textbf{Solution 3}

Use constructive counting, as per above. WLOG, place the red blocks first. There are 11 ways to place them with distance 0, 9 ways them to place with distance 2, so on, so the way to place red blocks is $11+9+7+5+3+1=36$. Then place any other block similarly, with $25$ ways (basic counting). You get then $6!^2$ ways to place the blocks evenly, and $12!/64$ ways to place the blocks in any way, so you get $\frac{16}{231}=247$ by simplifying.

 -drag00n

 





\textbf{Solution 4}

We can divide the $12$ positions into odd and even positions. Each color needs one block in an odd position and one block in an even position.

 WLOG, we place the first block of the first pair into an odd position. This leaves $6$ even positions out of the $11$ remaining positions for the second block of the first pair. This results in a probability of $\frac{6}{11}$ for the second block to fall into an even position.

 We can now place the first block of the second pair into another odd position, leaving $5$ even positions out of the $9$ remaining positions for the second block of the second pair.

 Continuing this pattern for the other $4$ pairs results in the product $\frac{6\cdot5\cdot4\cdot3\cdot2\cdot1}{11\cdot9\cdot7\cdot5\cdot3\cdot1}=\frac{16}{231}$. Thus, our answer is $16+231=\boxed{247}$.

 





\textbf{Solution 5}

There are a total of $\frac{12!}{(2!)^{6}}$ ways to arrange the blocks because there are $12!$ total ways to arrange the blocks but we're overcounting $2!$ ways for each of the $6$ blocks. Now, to better visualize how we can arrange these blocks, consider the first slot. There are $6$ colors to choose from. WLOG, say it's red. Now, we can easily see that the next and last red block can be put into a total of 6 slots to make sure that the number of slots between the first slot and the next slot is even. To better visualize this, I have drawn a chart type figure:

 RR\_ \_ \_ \_ \_ \_ \_ \_ \_ \_R \_ \_ R \_ \_ \_ \_ \_ \_ \_ \_ ...From this, you can get the idea. Counting up all these cases, it is obvious that there are $6$ ways to put the next and last red block. We continue by placing the next block into the next open slot. There are $5$ ways to choose the next color and similarly, only $5$ ways to put the next and last colored block to maintain an even number of slots. We get the idea. It seems that there are $(6!)^{2}$ ways to put the blocks. We compute $\frac{(6!)^{2} \cdot (2!)^{6}}{12!} = \frac{16}{231}$. This yields an answer of $16 + 231 = \boxed{247}$.

 




	\item (\textit{AIME I 2022}) Three spheres with radii $11,$ $13,$ and $19$ are mutually externally tangent. A plane intersects the spheres in three congruent circles centered at $A,$ $B,$ and $C,$ respectively, and the centers of the spheres all lie on the same side of this plane. Suppose that $AB^2 = 560.$ Find $AC^2.$


 
 \textbf{Answer: }756



\textbf{Solution 1}

This solution refers to the \textbf{Diagram} section.

 We let $\ell$ be the plane that passes through the spheres and $O_A$ and $O_B$ be the centers of the spheres with radii $11$ and $13$. We take a cross-section that contains $A$ and $B$, which contains these two spheres but not the third, as shown below:
                    \begin{figure}[H]
                        \centering
                        \includegraphics[width=0.5\textwidth]{Solution/AIME_2022_I_10_Solution_Figure_1.png}
                    \end{figure}
                    Because the plane cuts out congruent circles, they have the same radius and from the given information, $AB = \sqrt{560}$. Since $ABO_BO_A$ is a trapezoid, we can drop an altitude from $O_A$ to $BO_B$ to create a rectangle and triangle to use Pythagorean theorem. We know that the length of the altitude is $\sqrt{560}$ and let the distance from $O_B$ to $D$ be $x$. Then we have $x^2 = 576-560 \implies x = 4$.

 We have $AO_A = BD$ because of the rectangle, so $\sqrt{11^2-r^2} = \sqrt{13^2-r^2}-4$.Squaring, we have $121-r^2 = 169-r^2 + 16 - 8 \cdot \sqrt{169-r^2}$.Subtracting, we get $8 \cdot \sqrt{169-r^2} = 64 \implies \sqrt{169-r^2} = 8 \implies 169-r^2 = 64 \implies r^2 = 105$.We also notice that since we had $\sqrt{169-r^2} = 8$ means that $BO_B = 8$ and since we know that $x = 4$, $AO_A = 4$.

 We take a cross-section that contains $A$ and $C$, which contains these two spheres but not the third, as shown below:
                    \begin{figure}[H]
                        \centering
                        \includegraphics[width=0.5\textwidth]{Solution/AIME_2022_I_10_Solution_Figure_2.png}
                    \end{figure}
                    We have $CO_C = \sqrt{19^2-r^2} = \sqrt{361 - 105} = \sqrt{256} = 16$. Since $AO_A = 4$, we have $EO_C = 16-4 = 12$. Using Pythagorean theorem, $O_AE = \sqrt{30^2 - 12^2} = \sqrt{900-144} = \sqrt{756}$. Therefore, $O_AE^2 = AC^2  = \boxed{756}$.

 





\textbf{Solution 2}

Let the distance between the center of the sphere to the center of those circular intersections as $a,b,c$ separately.

 According to the problem, we have $a^2-11^2=b^2-13^2=c^2-19^2; (11+13)^2-(b-a)^2=560.$ After solving we have $b-a=4,$ plug this back to $11^2-a^2=13^2-b^2,$ we have $a=4, b=8,$ and $c=16.$

 The desired value is $(11+19)^2-(16-4)^2=\boxed{756}.$

 





\textbf{Solution 3}

Denote by $r$ the radius of three congruent circles formed by the cutting plane.Denote by $O_A$, $O_B$, $O_C$ the centers of three spheres that intersect the plane to get circles centered at $A$, $B$, $C$, respectively.

 Because three spheres are mutually tangent, $O_A O_B = 11 + 13 = 24$, $O_A O_C = 11 + 19 = 30$.

 We have $O_A A^2 = 11^2 - r^2$, $O_B B^2 = 13^2 - r^2$, $O_C C^2 = 19^2 - r^2$.

 Because $O_A A$ and $O_B B$ are perpendicular to the plane, $O_A AB O_B$ is a right trapezoid, with $\angle O_A A B = \angle O_B BA = 90^\circ$.

 Hence,
 \begin{align*} O_B B - O_A A & = \sqrt{O_A O_B^2 - AB^2} \\ & = 4 . \hspace{1cm} (1) \end{align*}


 Recall that
 \begin{align*} O_B B^2 - O_A A^2 & = \left( 13^2 - r^2 \right) - \left( 11^2 - r^2 \right) \\ & = 48 . \hspace{1cm} (2) \end{align*}


 Hence, taking $\frac{(2)}{(1)}$, we get
 \[ O_B B + O_A A = 12 . \hspace{1cm} (3) \]

 Solving (1) and (3), we get $O_B B = 8$ and $O_A A = 4$.

 Thus, $r^2 = 11^2 - O_A A^2 = 105$.

 Thus, $O_C C = \sqrt{19^2 - r^2} = 16$.

 Because $O_A A$ and $O_C C$ are perpendicular to the plane, $O_A AC O_C$ is a right trapezoid, with $\angle O_A A C = \angle O_C CA = 90^\circ$.

 Therefore,
 \begin{align*} AC^2 & = O_A O_C^2 - \left( O_C C - O_A A \right)^2 \\ & = \boxed{756}. \end{align*}


 $\textbf{FINAL NOTE:}$ In our solution, we do not use the condition that spheres $A$ and $B$ are externally tangent. This condition is redundant in solving this problem. 

 $\textbf{MORE FINAL NOTE:}$ the above note is incorrect because that condition was used at the start when claiming $O_AO_B=24$. Perhaps the note is referring to spheres $B$ and $C$.

 




	\item (\textit{AIME I 2022}) Let $ABCD$ be a parallelogram with $\angle BAD < 90^\circ.$ A circle tangent to sides $\overline{DA},$ $\overline{AB},$ and $\overline{BC}$ intersects diagonal $\overline{AC}$ at points $P$ and $Q$ with $AP < AQ,$ as shown. Suppose that $AP=3,$ $PQ=9,$ and $QC=16.$ Then the area of $ABCD$ can be expressed in the form $m\sqrt{n},$ where $m$ and $n$ are positive integers, and $n$ is not divisible by the square of any prime. Find $m+n.$


 
                    \begin{figure}[H]
                        \centering
                        \includegraphics[width=0.5\textwidth]{Problem/AIME_2022_I_11_Problem_Figure_1.png}
                    \end{figure}
                    


 
 \textbf{Answer: }150



\textbf{Solution 1 (No trig)}

Let's redraw the diagram, but extend some helpful lines.

 
                    \begin{figure}[H]
                        \centering
                        \includegraphics[width=0.5\textwidth]{Solution/AIME_2022_I_11_Solution_Figure_1.png}
                    \end{figure}
                    

 We obviously see that we must use power of a point since they've given us lengths in a circle and there are intersection points. Let $T_1, T_2, T_3$ be our tangents from the circle to the parallelogram. By the secant power of a point, the power of $A = 3 \cdot (3+9) = 36$. Then $AT_2  = AT_3 = \sqrt{36} = 6$. Similarly, the power of $C = 16 \cdot (16+9) = 400$ and  $CT_1 = \sqrt{400} = 20$. We let $BT_3 = BT_1 = x$ and label the diagram accordingly.

 Notice that because $BC = AD, 20+x = 6+DT_2 \implies DT_2 = 14+x$. Let $O$ be the center of the circle. Since $OT_1$ and $OT_2$ intersect $BC$ and $AD$, respectively, at right angles, we have $T_2T_1CD$ is a right-angled trapezoid and more importantly, the diameter of the circle is the height of the triangle. Therefore, we can drop an altitude from $D$ to $BC$ and $C$ to $AD$, and both are equal to $2r$. Since $T_1E = T_2D$, $20 - CE = 14+x \implies CE = 6-x$. Since $CE = DF, DF = 6-x$ and $AF = 6+14+x+6-x = 26$. We can now use Pythagorean theorem on $\triangle ACF$; we have $26^2 + (2r)^2 = (3+9+16)^2 \implies 4r^2 = 784-676 \implies 4r^2 = 108 \implies 2r = 6\sqrt{3}$ and $r^2 = 27$.

 We know that $CD = 6+x$ because $ABCD$ is a parallelogram. Using Pythagorean theorem on $\triangle CDF$, $(6+x)^2 = (6-x)^2 + 108 \implies (6+x)^2-(6-x)^2 = 108 \implies 12 \cdot 2x = 108 \implies 2x = 9 \implies x = \frac{9}{2}$. Therefore, base $BC = 20 + \frac{9}{2} = \frac{49}{2}$. Thus the area of the parallelogram is the base times the height, which is $\frac{49}{2} \cdot 6\sqrt{3} = 147\sqrt{3}$ and the answer is $\boxed{150}$

 





\textbf{Solution 2}

Let the circle tangent to $BC,AD,AB$ at $P,Q,M$ separately, denote that $\angle{ABC}=\angle{D}=\alpha$

 Using POP, it is very clear that $PC=20,AQ=AM=6$, let $BM=BP=x,QD=14+x$, using LOC in $\triangle{ABP}$,$x^2+(x+6)^2-2x(x+6)\cos\alpha=36+PQ^2$, similarly, use LOC in $\triangle{DQC}$, getting that $(14+x)^2+(6+x)^2-2(6+x)(14+x)\cos\alpha=400+PQ^2$. We use the second equation to minus the first equation, getting that $28x+196-(2x+12)\times14\times\cos\alpha=364$, we can get $\cos\alpha=\frac{2x-12}{2x+12}$.

 Now applying LOC in $\triangle{ADC}$, getting $(6+x)^2+(20+x)^2-2(6+x)\times(20+x)\times\frac{2x-12}{2x+12}=(3+9+16)^2$, solving this equation to get $x=\frac{9}{2}$, then $\cos\alpha=-\frac{1}{7}$, $\sin\alpha=\frac{4\sqrt{3}}{7}$, the area is $\frac{21}{2}\cdot\frac{49}{2}\cdot\frac{4\sqrt{3}}{7}=147\sqrt{3}$ leads to $\boxed{150}$

 





\textbf{Solution 3}

Denote by $O$ the center of the circle. Denote by $r$ the radius of the circle.Denote by $E$, $F$, $G$ the points that the circle meets $AB$, $CD$, $AD$ at, respectively.

 Because the circle is tangent to $AD$, $CB$, $AB$, $OE = OF = OG = r$, $OE \perp AD$, $OF \perp CB$, $OG \perp AB$.

 Because $AD \parallel CB$, $E$, $O$, $F$ are collinear. 

 Following from the power of a point, $AG^2 = AE^2 = AP \cdot AQ$. Hence, $AG = AE = 6$.

 Following from the power of a point, $CF^2 = CQ \cdot CP$. Hence, $CF = 20$.

 Denote $BG = x$. Because $DG$ and $DF$ are tangents to the circle, $BF = x$.

 Because $AEFB$ is a right trapezoid, $AB^2 = EF^2 + \left( AE - BF \right)^2$.Hence, $\left( 6 + x \right)^2 = 4 r^2 + \left( 6 - x \right)^2$.This can be simplified as 
 \[ 6 x = r^2 . \hspace{1cm} (1) \]

 In $\triangle ACB$, by applying the law of cosines, we have
 \begin{align*} AC^2 & = AB^2 + CB^2 - 2 AB \cdot CB \cos B \\ & = AB^2 + CB^2 + 2 AB \cdot CB \cos A \\ & = AB^2 + CB^2 + 2 AB \cdot CB \cdot \frac{AE - BF}{AB} \\ & = AB^2 + CB^2 + 2 CB \left( AE - BF \right) \\ & = \left( 6 + x \right)^2 + \left( 20 + x \right)^2 + 2 \left( 20 + x \right) \left( 6 - x \right) \\ & = 24 x + 676 . \end{align*}


 Because $AC = AP + PQ + QC = 28$, we get $x = \frac{9}{2}$.Plugging this into Equation (1), we get $r = 3 \sqrt{3}$.

 Therefore,
 \begin{align*} {\textrm Area} \ ABCD & = CB \cdot EF \\ & = \left( 20 + x \right) \cdot 2r \\ & = 147 \sqrt{3} . \end{align*}


 Therefore, the answer is $147 + 3 = \boxed{\textbf{(150) }}$.

 





\textbf{Solution 4}

Let $\omega$ be the circle, let $r$ be the radius of $\omega$, and let the points at which $\omega$ is tangent to $AB$, $BC$, and $AD$ be $X$, $Y$, and $Z$, respectively. Note that PoP on $A$ and $C$ with respect to $\omega$ yields $AX=6$ and $CY=20$. We can compute the area of $ABC$ in two ways:

 1. By the half-base-height formula, $[ABC]=r(20+BX)$.

 2. We can drop altitudes from the center $O$ of $\omega$ to $AB$, $BC$, and $AC$, which have lengths $r$, $r$, and $\sqrt{r^2-\frac{81}{4}}$. Thus, $[ABC]=[OAB]+[OBC]+[OAC]=r(BX+13)+14\sqrt{r^2-\frac{81}{4}}$.

 Equating the two expressions for $[ABC]$ and solving for $r$ yields $r=3\sqrt{3}$. 

 Let $BX=BY=a$. By the Parallelogram Law, $(a+6)^2+(a+20)^2=38^2$. Solving for $a$ yields $a=9/2$. Thus, $[ABCD]=2[ABC]=2r(20+a)=147\sqrt{3}$, for a final answer of $\boxed{150}$.

 





\textbf{Solution 5}


                    \begin{figure}[H]
                        \centering
                        \includegraphics[width=0.5\textwidth]{Solution/AIME_2022_I_11_Solution_Figure_2.png}
                    \end{figure}
                    Let $\omega$ be the circle, let $r$ be the radius of $\omega$, and let the points at which $\omega$ is tangent to $AB$, $BC$, and $AD$ be $H$, $K$, and $T$, respectively. PoP on $A$ and $C$ with respect to $\omega$ yields 
 \[AT=6, CK=20.\]

 Let $TG = AC, CG||AT.$

 In  $\triangle KGT$ $KT \perp BC,$$KT = \sqrt{GT^2 – (KC + AT)^2} = 6 \sqrt{3}=2r.$ 

 $\angle AOB = 90^{\circ}, OH \perp AB,  OH = r = \frac{KT}{2},$
 \[OH^2 = AH \cdot BH \implies  BH = \frac {9}{2}.\]

 Area is 
 \[(BK + KC) \cdot KT = (BH + KC) \cdot 2r = \frac{49}{2} \cdot 6\sqrt{3} = 147 \sqrt{3} \implies 147+3 = \boxed{\textbf{150}}.\]

 \textbf{vladimir.shelomovskii@gmail.com, vvsss}

 





\textbf{Solution 6 (Short and Sweet)}

Let $O$ be the center of the circle. Let points $M, N$ and $L$ be the tangent points of lines $BC, AD$ and $AB$ respectively to the circle. By Power of a Point, $({MC})^2=16\cdot{25} \Longrightarrow MC=20$. Similarly, $({AL})^2=3\cdot{12} \Longrightarrow AL=6$. Notice that $AL=AN=6$ since quadrilateral $LONA$ is symmetrical. Let $AC$ intersect $MN$ at $I$. Then, $\bigtriangleup{IMC}$ is similar to $\bigtriangleup{AIN}$. Therefore, $\frac{CI}{MC}=\frac{AI}{AN}$. Let the length of $PI=l$, then $\frac{25-l}{20}=\frac{3+l}{6}$. Solving we get $l=\frac{45}{13}$. Doing the Pythagorean theorem on triangles $IMC$ and $AIN$ for sides $MI$ and $IN$ respectively, we obtain the equation $\sqrt{(\frac{280}{13})^2-400} +\sqrt{(\frac{84}{13})^2-36}=MN=2r_1$ where $r_1$ denotes the radius of the circle. Solving, we get $MN=6\sqrt{3}$. Additionally, quadrilateral $OLBM$ is symmetrical so $OL=OM$. Let $OL=OM=x$ and extend a perpendicular foot from $B$ to $AD$ and call it $R$. Then, $\bigtriangleup{ABR}$ is right with $AR=6-x$, $AB=6+x$, and $RB=2r_1=MN=6\sqrt{3}$. Taking the difference of squares, we get $108=24x \Longrightarrow x=\frac{9}{2}$. The area of $ABCD$ is $MN\cdot{BC}=(20+x)\cdot{MN} \Longrightarrow \frac{49}{2}\cdot{6\sqrt{3}}=147\sqrt{3}$. Therefore, the answer is $147+3=\boxed{150}$

 





\textbf{Solution 7 (Intuitive, no trig, no weird auxiliary lines)}

Say that $BC$ is tangent to the circle at $X$ and $AD$ tangent at $Y$. Also, $H$ is the intersection of $XY$ (diameter) and $AC$ (diagonal). Then by power of a point with given info on $A$ and $C$ we get that $AY=6$ and $CX=20$. Note that $HAY \sim HCX$, and since $\frac{AY}{CX}=\frac{3}{10}$ we note that 
 \[\frac{AH}{CH} = \frac{AP+PH}{CQ+QH} = \frac{3+PH}{16+QH} =\frac{AY}{CX}=\frac{3}{10}\]. Since $PH+HQ=9$, we get that $PH=\frac{45}{13}$ and $QH=\frac{72}{13}$. This is the length information within the circle.The same triangle similarity also means that $\frac{YH}{XH}=\frac{3}{10}$, so if the radius of the circle is $r$ then we have $XH=\frac{20}{13}r$ and $YH = \frac{6}{13}r$. By power of a point on H, we can figure out $r$:
 \[XH\cdot YH = PH \cdot QG\]
 \[\frac{20}{13}r \cdot \frac{6}{13}r = \frac{45}{13} \cdot \frac{72}{13}\]and we get that $r = 3 \sqrt 3$. Thus, we have that the height of the parallelogram is $2r=6 \sqrt 3$ and we want to find $BC$. If $AB$ is tangent to the circle at $E$, then set $a = BX = BE$. Using pythagorean theorem, $AO^2+BO^2=AB^2$ and we can plug in diagram values: 
 \[(AY^2+OY^2)+(BX^2+OX)^2=AB^2\] 
 \[(6^2+(3 \sqrt 3)^2) + (a^2+(3 \sqrt 3)^2)=(a+6)^2.\] Solving, we get $a=\frac{9}{2}$ Finally, we have $[ABCD]=XY \cdot BC = 6 \sqrt 3 \cdot (20+\frac{9}{2}) \rightarrow \boxed{150}$

 





\textbf{Solution 8 (Ptolemy's Theorem + Power of Point + Pythagorean Theorem)}

Let $E$, $F$, $G$ be the circle's point of tangency with sides $AD$, $AB$, and $BC$, respectively. Let $O$ be the center of the inscribed circle.

 By Power of a Point, $AE^2 = AP \cdot AQ = 3(3+9) = 36$, so $AE = 6$. Similarly, $GC^2 = CQ \cdot CP = 16(16+9) = 400$, so $GC = 20$.

 Construct $GE$, and let $I$ be the point of intersection of $GE$ and $AC$. $GE \perp BC$ and $GE \perp AD$. By AA, $\triangle IGC \sim \triangle IEA$, and we have $\frac{AI}{IC} = \frac{AE}{GC} = \frac{3}{10}$. We also know $AI + IC = AC = 28$, so $AI = \frac{84}{13}$ and $IC = \frac{280}{13}$.

 Using Pythagorean Theorem on $\triangle IEA$ and $\triangle CIG$, we find that $EI = \frac{18\sqrt{3}}{13}$ and $IG = \frac{60\sqrt{3}}{13}$. Thus, $GE = EI + IG = 6\sqrt{3}$, and the radius of the circle is $3\sqrt{3}$.

 Construct $EF$, $FG$. $\angle AFO = \angle AEO = 90^{\circ}$, so $AEOF$ is cyclic. Similarly, $BFOG$ is cyclic.

 Now, we attempt to set up Ptolemy. Using Pythagorean Theorem on $\triangle AEO$, we find that $AO = 3\sqrt{7}$. By Ptolemy's Theorem, $(AE)(FO) + (AF)(EO) = (AO)(FE)$, from which we have $(6)(3\sqrt{3}) + (6)(3\sqrt{3}) = (3\sqrt{7})(FE)$ and $FE = 12\frac{\sqrt{3}}{\sqrt{7}}$. From Thales' Circle, $\triangle FGE$ is a right triangle, and $EF^2 + FG^2 = GE^2$, so $FG = \frac{18}{\sqrt{7}}$.

 Set $BF = BG = s$. $BO = \sqrt{s^2 + (3\sqrt{3})^2} = \sqrt{s^2+27}$, so by Ptolemy's Theorem on $BFOG$, we have

 
 \[(BF)(GO) + (BG)(FO) = (FG)(BO)\]
 \[(3\sqrt{3})(s) + (3\sqrt{3})(s) = (\frac{18}{\sqrt{7}})(\sqrt{s^2+27})\]Solving yields $s = \frac{9}{2}$.

 We know that $BC = BG + GC = 20 + \frac{9}{2} = \frac{49}{2}$, so the area of $ABCD = (\frac{49}{2})(6\sqrt{3}) = 147\sqrt{3}$. The requested answer is $147 + 3 = \boxed{150}$.

 





\textbf{Solution 9}

Let points $E$, $F$, $G$ be the points where lines $BC$, $AB$, and $AD$ are tangent to the circle respectively. Then extend line $BC$ to point $H$ such that $AH$ is perpendicular to $BC$.

 According to Power of a Point, $CQ\cdot CP=CE^2\rightarrow 16\cdot 25=CE^2\rightarrow CE=20$.

 Similarly, $AP\cdot AQ=AF^2\rightarrow 3\cdot 12=AF^2\rightarrow AF=6$.

 (Note that $AG=AF=6$ because they are intersecting tangents from the same circle. Furthermore, $HE=AG=6$ due to a simple upwards translation.)

 Therefore, $HC=HE+EC=6+20=26$, and we already know that $AC=28$, meaning that we can use the Pythagorean Theorem on $\Delta AHC$ to obtain: $AH=6\sqrt{3}$.

 Let $HB=k$. Then $BE=6-HB=6-k$. Since they are intersecting tangents from the same circle, $BF=BE=6-k$.

 Therefore, $AB=AF+BF=6+6-k=12-k$. We have all the side lengths of $\Delta ABH$, so applying the Pythagorean Theorem: $(6\sqrt{3})^2+k^2=(12-k)^2\rightarrow k=\frac{3}{2}$.

 Thus, $BC=BE+EC=6-k+20=26-k=\frac{49}{2}$.

 $[ABCD]=AH\cdot BC=6\sqrt{3}\cdot \frac{49}{2}=147\sqrt{3}$, so the answer is $\boxed{150}$.

 




	\item (\textit{AIME I 2022}) For any finite set $X,$ let $|X|$ denote the number of elements in $X.$ Define 
 \[S_n = \sum |A \cap B|,\] where the sum is taken over all ordered pairs $(A,B)$ such that $A$ and $B$ are subsets of $\{1,2,3,\ldots,n\}$ with $|A|=|B|.$ For example, $S_2 = 4$ because the sum is taken over the pairs of subsets 
 \[(A,B) \in \left\{(\emptyset,\emptyset),(\{1\},\{1\}),(\{1\},\{2\}),(\{2\},\{1\}),(\{2\},\{2\}),(\{1,2\},\{1,2\})\right\},\] giving $S_2 = 0+1+0+0+1+2=4.$ Let $\frac{S_{2022}}{S_{2021}} = \frac{p}{q},$ where $p$ and $q$ are relatively prime positive integers. Find the remainder when $p+q$ is divided by $1000.$


 
 \textbf{Answer: }245



\textbf{Solution 1 (Easy to Understand)}

Let's try out for small values of $n$ to get a feel for the problem. When $n=1, S_n$ is obviously $1$. The problem states that for $n=2, S_n$ is $4$. Let's try it out for $n=3$. 

 Let's perform casework on the number of elements in $A, B$. 

 $\textbf{Case 1:} |A| = |B| = 1$

 In this case, the only possible equivalencies will be if they are the exact same element, which happens $3$ times.

 $\textbf{Case 2:} |A| = |B| = 2$

 In this case, if they share both elements, which happens $3$ times, we will get $2$ for each time, and if they share only one element, which also happens $6$ times, we will get $1$ for each time, for a total of $12$ for this case.

 $\textbf{Case 3:} |A| = |B| = 3$

 In this case, the only possible scenario is that they both are the set $\{1,2,3\}$, and we have $3$ for this case.

 
In total, $S_3 = 18$.

 Now notice, the number of intersections by each element $1 \ldots 3$, or in general, $1 \ldots n$ is equal for each element because of symmetry - each element when $n=3$ adds $6$ to the answer. Notice that $6 = \binom{4}{2}$ - let's prove that $S_n = n \cdot \binom{2n-2}{n-1}$ (note that you can assume this and answer the problem if you're running short on time in the real test).

 Let's analyze the element $k$ - to find a general solution, we must count the number of these subsets that $k$ appears in. For $k$ to be in both $A$ and $B$, we need both sets to contain $k$ and another subset of $1$ through $n$ not including $k$. ($A = \{k\} \cup A'| A' \subset \{1,2,\ldots,n\} \land A' \not \subset \{k\}$ and $B = \{k\} \cup B'| B' \subset \{1,2,\ldots,n\} \land B' \not \subset \{k\}$)

 For any $0\leq l \leq n-1$ that is the size of both $A'$ and $B'$, the number of ways to choose the subsets $A'$ and $B'$ is $\binom{n-1}{l}$ for both subsets, so the total number of ways to choose the subsets are $\binom{n-1}{l}^2$.Now we sum this over all possible $l$'s to find the total number of ways to form sets $A$ and $B$ that contain $k$. This is equal to $\sum_{l=0}^{n-1} \binom{n-1}{l}^2$. This is a simplification of Vandermonde's identity, which states that $\sum_{k=0}^{r} \binom{m}{k} \cdot \binom{n}{r-k} = \binom{m+n}{r}$. Here, $m$, $n$ and $r$ are all $n-1$, so this sum is equal to $\binom{2n-2}{n-1}$. Finally, since we are iterating over all $k$'s for $n$ values of $k$, we have $S_n = n \cdot \binom{2n-2}{n-1}$, proving our claim.

 We now plug in $S_n$ to the expression we want to find. This turns out to be $\frac{2022 \cdot \binom{4042}{2021}}{2021 \cdot \binom{4040}{2020}}$. Expanding produces $\frac{2022 \cdot 4042!\cdot 2020! \cdot 2020!}{2021 \cdot 4040! \cdot 2021! \cdot 2021!}$.

 After cancellation, we have 
 \[\frac{2022 \cdot 4042 \cdot 4041}{2021 \cdot 2021 \cdot 2021} \implies \frac{4044\cdot 4041}{2021 \cdot 2021}\]

 $4044$ and $4041$ don't have any common factors with $2021$, so we're done with the simplification. We want to find $4044 \cdot 4041 + 2021^2 \pmod{1000} \equiv 44 \cdot 41 + 21^2 \pmod{1000} \equiv 1804+441 \pmod{1000} \equiv 2245 \pmod{1000} \equiv \boxed{245}$

 





\textbf{Solution 2 (Linearity of Expectation)}

We take cases based on the number of values in each of the subsets in the pair. Suppose we have $k$ elements in each of the subsets in a pair (for a total of n elements in the set). The expected number of elements in any random pair will be $n \cdot \frac{k}{n} \cdot \frac{k}{n}$ by linearity of expectation because for each of the $n$ elements, there is a $\frac{k}{n}$ probability that the element will be chosen. To find the sum over all such values, we multiply this quantity by  $\binom{n}{k}^2$. Summing, we get 
 \[\sum_{k=1}^{n} \frac{k^2}{n} \binom{n}{k}^2\] Notice that we can rewrite this as 
 \[\sum_{k=1}^{n} \frac{1}{n} \left(\frac{k \cdot n!}{(k)!(n - k)!}\right)^2 = \sum_{k=1}^{n} \frac{1}{n} n^2 \left(\frac{(n-1)!}{(k - 1)!(n - k)!}\right)^2 = n \sum_{k=1}^{n} \binom{n - 1}{k - 1}^2 = n \sum_{k=1}^{n} \binom{n - 1}{k - 1}\binom{n - 1}{n - k}\] We can simplify this using Vandermonde's identity to get $n \binom{2n - 2}{n - 1}$. Evaluating this for $2022$ and $2021$ gives 
 \[\frac{2022\binom{4042}{2021}}{2021\binom{4040}{2020}} = \frac{2022 \cdot 4042 \cdot 4041}{2021^3} = \frac{2022 \cdot 2 \cdot 4041}{2021^2}\] Evaluating the numerators and denominators mod $1000$ gives $804 + 441 = 1\boxed{245}$

 - pi\_is\_3.14

 





\textbf{Solution 3 (Rigorous)}

For each element $i$, denote $x_i = \left( x_{i, A}, x_{i, B} \right) \in \left\{ 0 , 1 \right\}^2$, where $x_{i, A} = \Bbb I \left\{ i \in A \right\}$ (resp. $x_{i, B} = \Bbb I \left\{ i \in B \right\}$).

 Denote $\Omega = \left\{ (x_1, \cdots , x_n): \sum_{i = 1}^n x_{i, A} = \sum_{i = 1}^n x_{i, B} \right\}$.

 Denote $\Omega_{-j} = \left\{ (x_1, \cdots , x_{j-1} , x_{j+1} , \cdots , x_n): \sum_{i \neq j} x_{i, A} = \sum_{i \neq j} x_{i, B} \right\}$.

 Hence,
 \begin{align*} S_n & = \sum_{(x_1, \cdots , x_n) \in \Omega} \sum_{i = 1}^n \Bbb I \left\{ x_{i, A} = x_{i, B} = 1 \right\} \\ & = \sum_{i = 1}^n \sum_{(x_1, \cdots , x_n) \in \Omega} \Bbb I \left\{ x_{i, A} = x_{i, B} = 1 \right\} \\ & = \sum_{i = 1}^n \sum_{(x_1, \cdots , x_{i-1} , x_{i+1} , \cdots , x_n) \in \Omega_{-i}} 1  \\ & = \sum_{i = 1}^n \sum_{j=0}^{n-1} \left( \binom{n-1}{j} \right)^2 \\ & = n \sum_{j=0}^{n-1} \left( \binom{n-1}{j} \right)^2 \\ & = n \sum_{j=0}^{n-1} \binom{n-1}{j} \binom{n-1}{n-1-j} \\ & = n \binom{2n-2}{n-1} . \end{align*}


 Therefore,
 \begin{align*} \frac{S_{2022}}{S_{2021}} & = \frac{2022 \binom{4042}{2021}}{2021 \binom{4040}{2020}} \\ & = \frac{4044 \cdot 4041}{2021^2} . \end{align*}


 This is in the lowest term.Therefore, modulo 1000,
 \begin{align*} p + q  & \equiv 4044 \cdot 4041 + 2021^2 \\ & \equiv 44 \cdot 41 + 21^2 \\ & \equiv \boxed{\textbf{(245) }} . \end{align*}


 





\textbf{Solution 4}

Let's ask what the contribution of an element $k\in \{1,2,\cdots,n\}$ is to the sum $S_n = \sum | A \cap B |.$

 The answer is given by the number of $(A,B)$ such that $|A|=|B|$ and $k \in A\cap B$, which is given by $\binom{2n-2}{n-1}$by the following construction: Write down 1 to $n$ except $k$ in a row. Do the same in a second row. Then choose $n-1$ numbers out of these $2n-2$ numbers. $k$ and the numbers chosen in the first row make up $A$. $k$ and the numbers not chosen in the second row make up $B$. This is a one-to-one correspondence between $(A,B)$ and the ways to choose $n-1$ numbers from $2n-2$ numbers.

 The contribution from all elements is therefore 
 \[S_n = n\binom{2n-2}{n-1}.\]For the rest please see Solution 1 or 2.

 NOTE: you may have a concern, "what if we chose the same element twice in a set?" Let us look at an example. Suppose we are partioning $S_4,$ $k=4,$ and we list out the rowstop: $1,2,3$bottom: $1,2,3$If we chose $1$ (top) $2$ (top) and $1$ (bottom), then we can distribute these numbers by sending the numbers in the top row that were chosen to set $A,$ which would be ${1,2,4},$ and set $B$ would be the numbers not chosen in the bottom row, which would be ${2,3,4}.$ To prove this rigourously, let $x$ be the number of elements in the top row chosen and $y$ be the number of elements chosen in the bottom row. Since we choose $n-1$ objects, $x+y=n-1.$ So the amount of elements in $A$ would be $1+x,$ and the amount of elements in $B$ would be $1+(n-1-y)=1+n-1-(n-1-x)=x+1.$

 





\textbf{Solution 5 (No strange notation, proof of all combo lemmas provided)}

Assume the subset $A$ of $\{1, 2, 3, \dots, n\}$ has $k$ elements, where $k$ is some integer from $0$ to $n$, inclusive. Then, there are $\tbinom{n}{k}$ ways to choose $A$. Now, we find the amount of ways to choose $B$, given $A$. Assume that $A$ and $B$ have $a$ elements in common, where $0 \le a \le k$. Then, $a$ of the elements in $B$ are also in $A$, so we can pick $\tbinom{k}{a}$ elements here, and the remaining $k-a$ elements in $B$ are $\underline{\text{not}}$ in $A$, so we can pick $\tbinom{n-k}{k-a}$ elements here. Therefore, there are $\tbinom{k}{a}\tbinom{n-k}{k-a}$ ways to select $B$ if $A$ and $B$ have $a$ elements in common.

 However, if $A$ and $B$ have $a$ elements in common, the ordered pair $(A, B)$ will contribute $a$ to the total sum. Thus, we multiply by $a$. Since the amount of ways to choose $B$ ranges over all $0 \le a \le k$, we seek the sum;

 $\binom{n}{k}\sum_{a=0}^{k} a\binom{k}{a}\binom{n-k}{k-a}$. 

 The $\binom{n}{k}$ on the outside of the sum represents the amount of ways to choose $A$. Now, we simplify the $a\tbinom{k}{a}$ term a little bit

 Note that $a\tbinom{k}{a}=a \cdot \frac{k!}{a!(k-a)!}=\frac{k!}{(a-1)!(k-a)!}=k \cdot \frac{(k-1)!}{(a-1)!(k-a)!}=k\binom{k-1}{a-1}$

 The reason we did this is because we want the $a$ in the combination expression, and not as a scale factor. Then, the sum simplifies to 

 $\binom{n}{k} \sum_{a=0}^{k} k\binom{k-1}{a-1}\binom{n-k}{k-a}$. We can take $k$ out of the sum (it's a constant in the sum), and the sum turns into

 $k\binom{n}{k}\sum_{a=0}^{k}\binom{k-1}{a-1}\binom{n-k}{k-a}$.

 Now, we use a combinatorial identity called Vandermode's Identity (This identity is well known, if you don't know it, learn it; it's useful for these types of problems). Consider $x+y$ people. We wish to form a committee with $r$ people. We can do this in $\tbinom{x+y}{r}$ ways. However, we can choose $k$ people from the group of $x$ people and $r-k$ people from the other $y$ people, for all $0 \le k \le r$. 

 This is done in $\sum_{k=0}^{r} \binom{x}{k} \binom{y}{r-k}$ ways, but this is also equal to $\binom{x+y}{r}$.

 In our desired sum, we have $x=k-1$, $y=n-k$, and $r=(a-1)+(k-a)=k-1$. Thus, the summation part is equal to $\binom{n-1}{k-1}$, and the total product is $k\binom{n}{k}\binom{n-1}{k-1}$.

 Now, we want to sum this over all possible $k$, for $0 \le k \le n$.

 This is $\sum_{k=0}^{n} k\binom{n}{k}\binom{n-1}{k-1}$. We can do something with the $k\binom{n}{k}$ term that was similar to something we did last time to obtain that this is also equal to $n\binom{n-1}{k-1}$.

 Then, the sum simplifies to $n\sum_{k=0}^{n} \binom{n-1}{k-1} \binom{n-1}{k-1}=n\sum_{k=0}^{n}\binom{n-1}{k-1}^{2}$.

 We can use Vandermode's again, with $x=y=n-1$ and $r=(k-1)+(k-1)=2k-2$. Then, the sum is equal to $n\binom{2n-2}{n-1}$.

 Now, the problem seeks $\frac{S_{2022}}{S_{2021}}$. This is equal to $\frac{2022\tbinom{4042}{2021}}{2021\tbinom{4040}{2020}}=\frac{\frac{2022 \cdot 4042!}{2021! \cdot 2021!}}{\frac{2021 \cdot 4040!}{2020! \cdot 2020!}}=\frac{2022 \cdot 4042! \cdot 2020! \cdot 2020!}{2021 \cdot 4040! \cdot 2021! \cdot 2021!}=\frac{2022 \cdot 4042 \cdot 4041}{2021 \cdot 2021 \cdot 2021}=\frac{2022 \cdot 2 \cdot 4041}{2021 \cdot 2021}$.

 We seek $2022 \cdot 2 \cdot 4041 + 2021 \cdot 2021 \pmod{1000}$. This is equivalent to $22 \cdot 82 + 21 \cdot 21 \equiv 804+441 \equiv \boxed{245} \pmod{1000}$, and we are done.

 




	\item (\textit{AIME I 2022}) Let $S$ be the set of all rational numbers that can be expressed as a repeating decimal in the form $0.\overline{abcd},$ where at least one of the digits $a,$ $b,$ $c,$ or $d$ is nonzero. Let $N$ be the number of distinct numerators obtained when numbers in $S$ are written as fractions in lowest terms. For example, both $4$ and $410$ are counted among the distinct numerators for numbers in $S$ because $0.\overline{3636} = \frac{4}{11}$ and $0.\overline{1230} = \frac{410}{3333}.$ Find the remainder when $N$ is divided by $1000.$


 
 \textbf{Answer: }392



\textbf{Solution 1}

$0.\overline{abcd}=\frac{abcd}{9999} = \frac{x}{y}$, $9999=9\times 11\times 101$.

 Then we need to find the number of positive integers $x$ that (with one of more $y$ such that $y|9999$) can meet the requirement $1 \leq {x}\cdot\frac{9999}{y} \leq 9999$.

 Make cases by factors of $x$. (A venn diagram of cases would be nice here.)

 
Case $A$:

 $3 \nmid x$ and $11 \nmid x$ and $101 \nmid x$, aka $\gcd (9999, x)=1$.

 Euler's totient function counts these: 
 \[\varphi \left(3^2 \cdot 11 \cdot 101 \right) = ((3-1)\cdot 3)(11-1)(101-1)= \bf{6000}\] values (but it's enough to note that it's a multiple of 1000 and thus does not contribute to the final answer)

 Note: You don't need to know this formula. The remaining cases essentially re-derive the same computation for other factors of $9999$. This case isn't actually different. 

 The remaining cases have $3$ (or $9$), $11$, and/or $101$ as factors of $abcd$, which cancel out part of $9999$.Note: Take care about when to use $3$ vs $9$.

 
Case $B$: $3|x$, but $11 \nmid x$ and $101 \nmid x$. 

 Then $abcd=9x$ to leave 3 uncancelled, and $x=3p$, so  $x \leq \frac{9999}{9} = 1111$, giving:

 $x \in 3 \cdot \{1, \dots \left\lfloor \frac{1111}{3}\right\rfloor\}$,

 $x \notin (3\cdot 11) \cdot \{1 \dots \left\lfloor \frac{1111}{3\cdot 11}\right\rfloor\}$,

 $x \notin (3 \cdot 101) \cdot \{1 \dots \left\lfloor \frac{1111}{3 \cdot 101}\right\rfloor\}$,

 for a subtotal of $\left\lfloor \frac{1111}{3}\right\rfloor -  (\left\lfloor\frac{1111}{3 \cdot 11}\right\rfloor + \left\lfloor\frac{1111}{3 \cdot 101}\right\rfloor ) = 370 - (33+3) = \bf{334}$ values.

 
Case $C$: $11|x$, but $3 \nmid x$ and $101 \nmid x$. 

 Much like previous case, $abcd$ is $11x$, so  $x \leq \frac{9999}{11} = 909$,  

 giving $\left\lfloor \frac{909}{11}\right\rfloor -  \left(\left\lfloor\frac{909}{11 \cdot 3}\right\rfloor + \left\lfloor\frac{909}{11 \cdot 101}\right\rfloor \right) = 82 - (27 + 0) = \bf{55}$ values.

 
Case $D$: $3|x$ and $11|x$ (so $33|x$), but $101 \nmid x$.  

 Here, $abcd$ is $99x$, so  $x \leq \frac{9999}{99} = 101$, 

 giving $\left\lfloor \frac{101}{33}\right\rfloor - \left\lfloor \frac{101}{33 \cdot 101}\right\rfloor = 3-0 = \bf{3}$ values.

 
Case $E$: $101|x$.

 Here, $abcd$ is $101x$, so  $x \leq \frac{9999}{101} = 99$,

 giving $\left\lfloor \frac{99}{101}\right\rfloor = \bf{0}$ values, so we don't need to account for multiples of $3$ and $11$. 

 To sum up, the answer is 
 \[6000+334+55+3+0\equiv\boxed{392} \pmod{1000}.\]

 \textbf{Clarification} In this context, when the solution says, "Then $abcd=9x$ to leave 3 uncancelled, and $x=3p$," it is a bit vague. The best way to clarify this is by this exact example - what is really meant is we need to divide by 9 first to achieve 1111, which has no multiple of 3; thus, given that the fraction x/y is the simplest form, x can be a multiple of 3.

 Similar explanations can be said when the solution divides 9999 by 11, 101, and uses that divided result in the PIE calculation rather than 9999.

 \textbf{mathboy282}

 


 \textbf{Sketch of Constuctive Counting Solution} If anyone wants to write this up, basically I thought to do:
 \[ \sum_{d\mid 9999}\varphi(d)=9999 \]and which is 9999 (identity), but then we must subtract all the numbers that show up multiple times in the numerator. We do this by doing casework on $n,$ like if $1111<n\leq3333$ we check how many n are coprime to 3333 in this range, giving the amount of n that show up multiple times in both $\frac{n}{3333}$ and $\frac{n}{9999}.$ repeat this bounding process and it should work.

 \href{https://MathProblemSolvingSkills.com}{https://MathProblemSolvingSkills.com}

 




	\item (\textit{AIME I 2022}) Given $\triangle ABC$ and a point $P$ on one of its sides, call line $\ell$ the $\textit{splitting line}$ of $\triangle ABC$ through $P$ if $\ell$ passes through $P$ and divides $\triangle ABC$ into two polygons of equal perimeter. Let $\triangle ABC$ be a triangle where $BC = 219$ and $AB$ and $AC$ are positive integers. Let $M$ and $N$ be the midpoints of $\overline{AB}$ and $\overline{AC},$ respectively, and suppose that the splitting lines of $\triangle ABC$ through $M$ and $N$ intersect at $30^\circ.$ Find the perimeter of $\triangle ABC.$


 
 \textbf{Answer: }459



\textbf{Solution (Geometry + Number Theory)}

Denote $BC = a$, $CA = b$, $AB = c$.

 Let the splitting line of $\triangle ABC$ through $M$ (resp. $N$) crosses $\triangle ABC$ at another point $X$ (resp. $Y$).

 WLOG, we assume $c \leq b$.

 $\textbf{Case 1}$: $a \leq c \leq b$.

 We extend segment $AB$ to $D$, such that $BD = a$.We extend segment $AC$ to $E$, such that $CE = a$.

 In this case, $X$ is the midpoint of $AE$, and $Y$ is the midpoint of $AD$.

 Because $M$ and $X$ are the midpoints of $AB$ and $AE$, respectively, $MX \parallel BE$.Because $N$ and $Y$ are the midpoints of $AC$ and $AD$, respectively, $NY \parallel CD$.

 Because $CB = CE$, $\angle CBE =\angle CEB = \frac{\angle ACB}{2}$.Because $BC = BD$, $\angle BCD = \angle BDC = \frac{\angle ABC}{2}$.

 Let $BE$ and $CD$ intersect at $O$.Because $MX \parallel BE$ and $NY \parallel CD$, the angle formed between lines $MX$ and $NY$ is congruent to $\angle BOD$. Hence, $\angle BOD = 30^\circ$ or $150^\circ$.

 We have
 \begin{align*} \angle BOD & = \angle CBE + \angle BCD \\ & = \frac{\angle ACB}{2} + \frac{\angle ABC}{2} \\ & = 90^\circ - \frac{\angle A}{2} . \end{align*}


 Hence, we must have $\angle BOD = 30^\circ$, not $150^\circ$.Hence, $\angle A = 120^\circ$.

 This implies $a > b$ and $a >c$. This contradicts the condition specified for this case.

 Therefore, this case is infeasible.

 $\textbf{Case 2}$: $c \leq a \leq b$.

 We extend segment $CB$ to $D$, such that $BD = c$.We extend segment $AC$ to $E$, such that $CE = a$.

 In this case, $X$ is the midpoint of $AE$, and $Y$ is the midpoint of $CD$.

 Because $M$ and $X$ are the midpoints of $AB$ and $AE$, respectively, $MX \parallel BE$.Because $N$ and $Y$ are the midpoints of $AC$ and $CD$, respectively, $NY \parallel AD$.

 Because $CB = CE$, $\angle CBE =\angle CEB = \frac{\angle ACB}{2}$.Because $BA = BD$, $\angle BAD = \angle BDA = \frac{\angle ABC}{2}$.

 Let $O$ be a point of $AC$, such that $BO \parallel AD$.Hence, $\angle OBC = \angle BDA = \frac{B}{2}$.

 Because $MX \parallel BE$ and $NY \parallel AD$ and $AD \parallel BO$, the angle formed between lines $MX$ and $NY$ is congruent to $\angle OBE$. Hence, $\angle OBE = 30^\circ$ or $150^\circ$.

 We have
 \begin{align*} \angle OBE & = \angle OBC + \angle CBE \\ & = \frac{\angle ABC}{2} + \frac{\angle ACB}{2} \\ & = 90^\circ - \frac{\angle A}{2} . \end{align*}


 Hence, we must have $\angle OBE = 30^\circ$, not $150^\circ$.Hence, $\angle A = 120^\circ$.

 This implies $a > b$ and $a >c$. This contradicts the condition specified for this case.

 Therefore, this case is infeasible.

 $\textbf{Case 3}$: $c \leq b \leq a$.

 We extend segment $CB$ to $D$, such that $BD = c$.We extend segment $BC$ to $E$, such that $CE = b$.

 In this case, $X$ is the midpoint of $BE$, and $Y$ is the midpoint of $CD$.

 Because $M$ and $X$ are the midpoints of $AB$ and $BE$, respectively, $MX \parallel AE$.Because $N$ and $Y$ are the midpoints of $AC$ and $CD$, respectively, $NY \parallel AD$.

 Because $CA = CE$, $\angle CAE =\angle CEB = \frac{\angle ACB}{2}$.Because $BA = BD$, $\angle BAD = \angle BDA = \frac{\angle ABC}{2}$.

 Because $MX \parallel AE$ and $NY \parallel AD$, the angle formed between lines $MX$ and $NY$ is congruent to $\angle DAE$. Hence, $\angle DAE = 30^\circ$ or $150^\circ$.

 We have
 \begin{align*} \angle DAE & = \angle BAD + \angle CAE + \angle BAC \\ & = \frac{\angle ABC}{2} + \frac{\angle ACB}{2} + \angle BAC \\ & = 90^\circ + \frac{\angle BAC}{2} . \end{align*}


 Hence, we must have $\angle OBE = 150^\circ$, not $30^\circ$.Hence, $\angle BAC = 120^\circ$.

 In $\triangle ABC$, by applying the law of cosines, we have
 \begin{align*} a^2 & = b^2 + c^2 - 2bc \cos \angle BAC\\ & =  b^2 + c^2 - 2bc \cos 120^\circ \\ & = b^2 + c^2 + bc . \end{align*}


 Because $a = 219$, we have
 \[ b^2 + c^2 + bc  = 219^2 . \]

 Now, we find integer solution(s) of this equation with $c \leq b$.

 Multiplying this equation by 4, we get
 \[ \left( 2 c + b \right)^2 + 3 b^2 = 438^2 . \hspace{1cm} (1) \]

 Denote $d = 2 c + b$. Because $c \leq b$, $b < d \leq 3 b$.

 Because $438^2 - 3 b^2 \equiv 0 \pmod{3}$, $d^2 \equiv 0 \pmod{3}$.Thus, $d \equiv 0 \pmod{3}$.This implies $d^2 \equiv 0 \pmod{9}$. 

 We also have $438^2 \equiv 0 \pmod{9}$.Hence, $3 b^2 \equiv 0 \pmod{9}$.This implies $b \equiv 0 \pmod{3}$.

 Denote $b = 3 p$ and $d = 3 q$. Hence, $p < q \leq 3 p$.Hence, Equation (1) can be written as
 \[ q^2 + 3 p^2 = 146^2 . \hspace{1cm} (2) \]

 Now, we solve this equation.

 First, we find an upper bound of $q$.

 We have $q^2 + 3 p^2 \geq q^2 + 3 \left( \frac{q}{3} \right)^2 = \frac{4 q^2}{3}$.Hence, $\frac{4 q^2}{3} \leq 146^2$.Hence, $q \leq 73 \sqrt{3} < 73 \cdot 1.8 = 131.4$.Because $q$ is an integer, we must have $q \leq 131$.

 Second, we find a lower bound of $q$.

 We have $q^2 + 3 p^2 < q^2 + 3 q^2 = 4 q^2$.Hence, $4 q^2 > 146^2$.Hence, $q > 73$.Because $q$ is an integer, we must have $q \geq 74$.

 Now, we find the integer solutions of $p$ and $q$ that satisfy Equation (2) with $74 \leq q \leq 131$.

 First, modulo 9, 
 \begin{align*} q^2 & \equiv 146^2 - 3 p^2 \\ & \equiv 4 - 3 \cdot ( 0 \mbox{ or } 1 ) \\ & \equiv 4 \mbox{ or } 1 . \end{align*}


 Hence $q \equiv \pm 1, \pm 2 \pmod{9}$.

 Second, modulo 5, 
 \begin{align*} q^2 & \equiv 146^2 - 3 p^2 \\ & \equiv 1 + 2 p^2 \\ & \equiv 1 + 2 \cdot ( 0 \mbox{ or } 1 \mbox{ or } -1 ) \\ & \equiv 1 \mbox{ or } 3 \mbox{ or } - 1 . \end{align*}


 Because $q^2 \equiv 0 \mbox{ or } 1 \mbox{ or } - 1$, we must have $q^2 \equiv 1 \mbox{ or } - 1$.Hence, $5 \nmid q$.

 Third, modulo 7,
 \begin{align*} q^2 & \equiv 146^2 - 3 p^2 \\ & \equiv 1 - 3 \cdot ( 0 \mbox{ or } 1 \mbox{ or } 5 \mbox{ or } 2 ) \\ & \equiv 1 \mbox{ or } 2 \mbox{ or } 3 \mbox{ or } 5 . \end{align*}


 Because $q^2 \equiv 0 \mbox{ or } 1 \mbox{ or } 2 \mbox{ or } 4 \pmod{ 7 }$, we must have $q^2 \equiv 1 \mbox{ or } 2 \pmod{7}$.Hence, $q \equiv 1, 3, 4, 6 \pmod{7}$.

 Given all conditions above, the possible $q$ are 74, 83, 88, 92, 97, 101, 106, 109, 116, 118, 127.

 By testing all these numbers, we find that the only solution is $q = 97$.This implies $p = 63$.

 Hence, $b = 3p = 189$ and $d = 3q = 291$.Hence, $c = \frac{d - b}{2} = 51$.

 Therefore, the perimeter of $\triangle ABC$ is $b + c + a = 189 + 51 + 219 = \boxed{\textbf{(459) }}$.

 


 





\textbf{Solution (Number Theory Part)}

We wish to solve the Diophantine equation $a^2+ab+b^2=3^2 \cdot 73^2$. It can be shown that $3|a$ and $3|b$, so we make the substitution $a=3x$ and $b=3y$ to obtain $x^2+xy+y^2=73^2$ as our new equation to solve for. 

 Notice that $r^2+r+1=(r-\omega)(r-{\omega}^2)$, where $\omega=e^{i\frac{2\pi}{3}}$. Thus,
 \[x^2+xy+y^2 = y^2((x/y)^2+(x/y)+1) = y^2 (\frac{x}{y}-\omega)(\frac{x}{y}-{\omega}^2) = (x-y\omega)(x-y{\omega}^2).\]

 Note that $8^2+1^2+8 \cdot 1=73$. Thus, $(8-\omega)(8-{\omega}^2)=73$. Squaring both sides yields 
 \begin{align} (8-\omega)^2(8-{\omega}^2)^2&=73^2\\ (63-17\omega)(63-17{\omega}^2)&=73^2. \end{align}Thus, by $(2)$, $(63, 17)$ is a solution to $x^2+xy+y^2=73^2$. This implies that $a=189$ and $b=51$, so our final answer is $189+51+219=\boxed{459}$.

 





\textbf{Solution(Visual geometry)}


                    \begin{figure}[H]
                        \centering
                        \includegraphics[width=0.5\textwidth]{Solution/AIME_2022_I_14_Solution_Figure_1.png}
                    \end{figure}
                    
                    \begin{figure}[H]
                        \centering
                        \includegraphics[width=0.5\textwidth]{Solution/AIME_2022_I_14_Solution_Figure_2.png}
                    \end{figure}
                    
                    \begin{figure}[H]
                        \centering
                        \includegraphics[width=0.5\textwidth]{Solution/AIME_2022_I_14_Solution_Figure_3.png}
                    \end{figure}
                    We look at  upper and middle  diagrams and get $\angle BAC = 120^\circ$.

 Next we use only the lower Diagram. Let  $I$ be incenter $\triangle ABC$, E be midpoint of biggest arc $\overset{\Large\frown} {BC}.$Then bisector $AI$ cross circumcircle $\triangle ABC$ at point $E$. Quadrilateral  $ABEC$ is cyclic, so 
 \[\angle BEC = 180^\circ - \angle ABC = 60^\circ \implies BE = CE = IE = BC.\]
 \[AE \cdot BC = AB \cdot CE + AC \cdot  BE \implies AE = AB + AC\]$\implies AI +EI = AB + AC, \hspace{10mm} AI = AB+ AC – BC$ is integer.
 \[AI = \frac {2AB \cdot AC \cdot  cos \angle CAI}{AB+AC + BC} =  \frac {AB \cdot AC}{AB+AC + BC} =\]
 \[= AB + AC – BC \implies AC^2 + AB^2 + AB \cdot AC = BC^2.\]A quick $(\mod9)$ check gives that $3\mid AC$ and $3\mid AB$.
 \[AI \le A_0I_0 = EA_0 – EI_0 = \frac{2 BC}{\sqrt{3}} – BC = \frac {2 - \sqrt{3}}{\sqrt{3}} BC = 33.88.\]Denote $a= \frac {BC}{3}= 73, b = \frac {AC}{3}, c = \frac {AB}{3}, l = \frac {AI}{3} \le 11.$

 We have equations in integers$\frac{bc}{a+b+c} = b + c – a = l \le 11.$

 The solution $(b > c)$ is
 \[b = \frac{a + l +\sqrt{a^2 – 6al – 3l^2}}{2}, c = \frac{a + l -\sqrt{a^2 – 6al – 3l^2}}{2}.\]Suppose,  $a^2 – 6al – 3l^2 = (a – 3l – t)^2 \implies \frac {12l^2}{t} + t + 6l= 2a = 146.$

 Now we check all possible $t = {2,3,4,6,12, ml}.$

 Case  $t = 2 \implies 6l^2 + 6l = 146 – 2 \implies l^2 + l = 24 \implies \O$

 Case  $t = 3 \implies 4l^2 + 6l = 146 – 3 = 143\implies \O.$

 Case $t = 4 \implies 3l^2 + 6l = 146 – 4 =142 \implies \O.$

 Case $t = 6 \implies 2l^2 + 6l = 146 – 6 = 140 \implies l = 7, b = 63, c = 17.$

 Case $t = 12 \implies l^2 + 6l = 146 – 12 = 134 \implies \O.$

 Case $t = ml \implies \frac{12l}{m} + 6l + ml = 146 \implies \frac{12}{m} + 6 + m = \frac{73 \cdot 2}{l}\implies \O.$

 \textbf{vladimir.shelomovskii@gmail.com, vvsss}

 





\textbf{Solution (The Geometry Part Using Menelaus)}


                    \begin{figure}[H]
                        \centering
                        \includegraphics[width=0.5\textwidth]{Solution/AIME_2022_I_14_Solution_Figure_4.png}
                    \end{figure}
                    

 Let the 2 splitting lines crossing $N$ and $M$ cross $BC$ at $G$ and $H$. Extend $AB$ and $NG$ so that the two lines intersect at $E$ and extend $AC$ and $MH$ so that the two lines intersect at $F$. Let $a$,$b$,$c$ be the corresponding side lengths for $\angle A$,$\angle B$,$\angle C$. According to the question, $MH$ and $NG$ creates two sections of $\triangle ABC$ of equal perimeter, we could list out the equations:

 $BG+c=a-BG$

 $CH+b=a-CH$

 $BG+GH+CH=a$

 This results in:

 $BG=\frac{a-c}{2}$

 $CH=\frac{a-b}{2}$

 $GH=\frac{b+c}{2}$

 Now, we can apply the Menelaus Theorem to $\triangle ABC$. Let $BE = x$:
 \[\frac{BE}{AE}\cdot\frac{AN}{CN}\cdot\frac{CG}{BG} = 1\]
 \[\frac{x}{x+c}\cdot1\cdot\frac{\frac{a+c}{2}}{\frac{a-c}{2}} = 1\]
 \[\frac{x}{x+c}\cdot\frac{a+c}{a-c} = 1\]
 \[\frac{x+c}{x}=\frac{a+c}{a-c}\]
 \[\frac{c}{x}=\frac{2c}{a-c}\]
 \[BE=x=\frac{a-c}{2}=BG\]Similarly, $CF=CH$.

 Therefore, $\angle BGE = \angle E$ and $\angle CHF = \angle F$. Because the splitting lines of $\triangle ABC$ through $M$ and $N$ intersect at $30^\circ$, $\angle NOF$ = $30^\circ$. $\angle BGE+\angle CHF=\angle BGE+\angle CHF=\angle NOF=30^\circ$.

 We conclude that $\angle A=360^\circ-\angle E-\angle F-\angle EOF(reflex)=360^\circ-30^\circ-210^\circ=120^\circ$

 Apply the Law of Cosines on $\triangle ABC$

 $c^2+a^2-2ac\cos 120^\circ=219^2$

 $a^2+c^2+ac=219^2=47961$

 To find the final answer, you would only need to solve this equation, which steps could be found in previous solutions.

 





\textbf{Solution (Coordinate Geometry)}

Let the two $splitting \text{ }line$s be $\overline{MD}$ and $\overline{NE}$. We can approach this problem using coordinates.

 \textbf{Case 1:} $D$ lies on $BC$, $E$ lies on $BC$.

 Let $a=BC,b=CA,c=AB,\angle CAB=\angle A,\angle ABC=\angle B,\angle BCA=\angle C$. We establish coordinate system with $A$ being the origin and $\overline{AB}$ being the horizonal $x$-axis. Here is our setup:\begin{equation*}\begin{aligned}&A=(0,0) \\&B=(c,0) \\&C=(b\cos{\angle A},b\sin{\angle A}) \\&\overrightarrow{BC}=(b\cos{\angle A}-c,b\sin{\angle A}) \\&D=B+\frac{a+b}{2a}\cdot\overrightarrow{BC} \\&E=B+\frac{a-c}{2a}\cdot\overrightarrow{BC} \\&M=(\frac{c}{2},0) \\&N=(\frac{b}{2}\cos{\angle A},\frac{b}{2}\sin{\angle A}) \\&\overrightarrow{MD}=D-M \\&\overrightarrow{NE}=E-N\end{aligned}\end{equation*}Computing vector $\overrightarrow{MD},\overrightarrow{NE}$ yields\begin{equation*}\begin{aligned}&\overrightarrow{MD}=\frac{b}{2a}\cdot((a+b)\cos{\angle A}-c,(a+b)\sin{\angle A})\\&\overrightarrow{NE}=\frac{c}{2a}\cdot(a+c-b\cos{\angle A},-b\sin{\angle A})\end{aligned}\end{equation*}Let $\vec{u}=((a+b)\cos{\angle A}-c,(a+b)\sin{\angle A}),\ \vec{v}=(a+c-b\cos{\angle A},-b\sin{\angle A})$. Then for Case 1 we have
 \[\overrightarrow{MD}=\frac{b}{2a}\vec{u},\qquad\overrightarrow{NE}=\frac{c}{2a}\vec{v}.\]Note that for the other cases (arising from different relative lengths of the sides), the vectors $\overrightarrow{MD}$ and $\overrightarrow{NE}$ are simply $\vec{u}$ and $\vec{v}$ multiplied by a constant factor. The cases are:

 \textbf{Case 2:} $D$ on $BC$, $E$ on $AB$.
 \[\overrightarrow{MD}=\frac{b}{2a}\,\vec{u},\qquad \overrightarrow{NE}=\frac{1}{2}\,\vec{v}.\]

 \textbf{Case 3:} $D$ on $AC$, $E$ on $BC$.
 \[\overrightarrow{MD}=\frac{1}{2}\,\vec{u},\qquad \overrightarrow{NE}=\frac{c}{2a}\,\vec{v}.\]

 \textbf{Case 4:} $D$ on $AC$, $E$ on $AB$.
 \[\overrightarrow{MD}=\frac{1}{2}\,\vec{u},\qquad \overrightarrow{NE}=\frac{1}{2}\,\vec{v}.\]

 In every case, the common factors cancel when computing the cosine of the intersecting angle $\theta$, because
 \[\cos{\theta}=\frac{\lvert\overrightarrow{MD}\cdot\overrightarrow{NE}\rvert} {\lvert\overrightarrow{MD}\rvert\,\lvert\overrightarrow{NE}\rvert}=\frac{\lvert\vec{u}\cdot\vec{v}\rvert} {\lvert\vec{u}\rvert\,\lvert\vec{v}\rvert}.\]
 \[\cos{\theta}=\frac{\lvert\vec{u}\cdot\vec{v}\rvert}{\lvert\vec{u}\rvert\lvert\vec{v}\rvert}\]We proceed by computing $\cos{\theta}$. The result of the computation of the dot product and the magnitude holds\begin{equation*}\begin{aligned}&\lvert\vec{u}\cdot\vec{v}\rvert=\lvert a(\cos{\angle A}-1)(a+b+c)\rvert\\&\lvert\vec{u}\rvert=\sqrt{2a(a+b-c\cos{\angle A})}\\&\lvert\vec{v}\rvert=\sqrt{2a(a+c-b\cos{\angle A})}\end{aligned}\end{equation*}Using the identity $\cos{\angle A}-1=-2\sin^{2}{\frac{\angle A}{2}}$, the dot product simplifies to
 \[\lvert\vec{u}\cdot\vec{v}\rvert=2a(a+b+c)\sin^{2}{\frac{\angle A}{2}}\]By the Law of Cosine, $\cos{\angle A}=\frac{b^{2}+c^{2}-a^{2}}{2bc}$, we see that
 \[a+b-c\cos{\angle A}=\frac{2ab+2b^2-b^2-c^2+a^2}{2b}=\frac{(a+b)^2-c^2}{2b}\]Substituting the half perimeter $s=\frac{a+b+c}{2}$, we have\begin{equation*}\begin{aligned}&\lvert\vec{u}\cdot\vec{v}\rvert=4as\sin^{2}{\frac{\angle A}{2}}\\&\lvert\vec{u}\rvert=\sqrt{\frac{4as(s-c)}{b}}\\&\lvert\vec{v}\rvert=\sqrt{\frac{4as(s-b)}{c}}\end{aligned}\end{equation*}Therefore, $\cos{\theta}$ simplifies to
 \[\cos{\theta}=\sin^{2}{\frac{\angle A}{2}}\cdot\sqrt{\frac{bc}{(s-b)(s-c)}}\]Note the well known identity\begin{align*}\frac{(s-b)(s-c)}{bc}&=\frac{(a+c-b)(a+b-c)}{4bc}\\&=\frac{a^2-(b-c)^2}{4bc}\\&=\frac{1-\frac{b^2+c^2-a^2}{2bc}}{2}\\&=\sin^{2}{\frac{\angle A}{2}}\end{align*}
We obtain the beautiful equation
 \[\cos{\theta}=\sin{\frac{\angle A}{2}}\]It is given that $\theta=30^{\circ}$, solving yields $\angle A=120^{\circ}$By the Law of Cosine, we wish to solve the Diophantine equation
 \[b^2+bc+c^2=219^2\]Solving it using modular arithmetic or Eisenstein integers, as shown in the previous solutions, we have $b=51,c=189$. This leads to the final answer $a+b+c=219+51+189=\boxed{459}$

 


 \textbf{Animated Video Solution} 




	\item (\textit{AIME I 2022}) Let $x,$ $y,$ and $z$ be positive real numbers satisfying the system of equations:

 \begin{align*} \sqrt{2x-xy} + \sqrt{2y-xy} &= 1 \\ \sqrt{2y-yz} + \sqrt{2z-yz} &= \sqrt2 \\ \sqrt{2z-zx} + \sqrt{2x-zx} &= \sqrt3. \end{align*}
 
Then $\left[ (1-x)(1-y)(1-z) \right]^2$ can be written as $\frac{m}{n},$ where $m$ and $n$ are relatively prime positive integers. Find $m+n.$


 
 \textbf{Answer: }033



\textbf{Solution 1 (geometric interpretation)}

First, let define a triangle with side lengths $\sqrt{2x}$, $\sqrt{2z}$, and $l$, with altitude from $l$'s equal to $\sqrt{xz}$. $l = \sqrt{2x - xz} + \sqrt{2z - xz}$, the left side of one equation in the problem.

 Let  $\theta$ be angle opposite the side with length $\sqrt{2x}$. Then the altitude has length $\sqrt{2z} \cdot \sin(\theta) = \sqrt{xz}$ and thus $\sin(\theta) = \sqrt{\frac{x}{2}}$, so $x=2\sin^2(\theta)$ and the side length $\sqrt{2x}$ is equal to $2\sin(\theta)$.

 We can symmetrically apply this to the two other equations/triangles. 

 By law of sines, we have $\frac{2\sin(\theta)}{\sin(\theta)} = 2R$, with $R=1$ as the circumradius, same for all 3 triangles. The circumcircle's central angle to a side is $2 \arcsin(l/2)$, so the 3 triangles' $l=1, \sqrt{2}, \sqrt{3}$,  have angles $120^{\circ}, 90^{\circ}, 60^{\circ}$, respectively.

 This means that by half angle arcs, we see that we have in some order, $x=2\sin^2(\alpha)$, $y=2\sin^2(\beta)$, and $z=2\sin^2(\gamma)$ (not necessarily this order, but here it does not matter due to symmetry), satisfying that $\alpha+\beta=180^{\circ}-\frac{120^{\circ}}{2}$, $\beta+\gamma=180^{\circ}-\frac{90^{\circ}}{2}$, and $\gamma+\alpha=180^{\circ}-\frac{60^{\circ}}{2}$. Solving, we get $\alpha=\frac{135^{\circ}}{2}$, $\beta=\frac{105^{\circ}}{2}$, and $\gamma=\frac{165^{\circ}}{2}$.

 We notice that 
 \[[(1-x)(1-y)(1-z)]^2=[\sin(2\alpha)\sin(2\beta)\sin(2\gamma)]^2=[\sin(135^{\circ})\sin(105^{\circ})\sin(165^{\circ})]^2\] 
 \[=\left(\frac{\sqrt{2}}{2} \cdot \frac{\sqrt{6}-\sqrt{2}}{4} \cdot \frac{\sqrt{6}+\sqrt{2}}{4}\right)^2 = \left(\frac{\sqrt{2}}{8}\right)^2=\frac{1}{32} \to \boxed{033}. \blacksquare\]

 - kevinmathz

 





\textbf{Solution 2 (Condensed Algebraic Trig)}

Let $\sin{A} = \sqrt{\frac{x}{2}}$, $\sin{B} = \sqrt{\frac{y}{2}}$, and $\sin{C} = \sqrt{\frac{z}{2}}$ for some $A,B,C \in (0^\circ,90^\circ)$; therefore, $\cos{A} = \sqrt{\frac{2-x}{2}}$, etc. Dividing all three given equations by $2$, we see that they are equivalent to\begin{align*}\sin(A+B) &= \sin(30^\circ) \\\sin(B+C) &= \sin(45^\circ) \\\sin(C+A) &= \sin(60^\circ),\end{align*}
which produces $A=\frac{45}{2}^\circ$, $B=\frac{15}{2}^\circ$, and $C=\frac{75}{2}^\circ$. Thus our answer is\begin{align*}[(1-x)(1-y)(1-z)]^2 &= [(1-2\sin^2 A)(1-2\sin^2 B)(1-2\sin^2 C)]^2 \\&= (\cos(2A)\cos(2B)\cos(2C))^2 \\&= \left(\frac{\sqrt{2}}{2} \cdot \frac{\sqrt{6}+\sqrt{2}}{2} \cdot \frac{\sqrt{6}-\sqrt{2}}{2}\right)^2 \\&= \frac{1}{32},\end{align*}
producing the answer of $\boxed{033}$.

 Remark. Note that liberties have been taken in mathematical rigor in the above solution, especially when dealing with angles $A,B,C$ and the inverse sine function. However, these details do not affect the overall solution and may be carefully checked during the contest if desired; thus we feel justified in leaving them out to emphasize the elegance of the trigonometric substitution.

 





\textbf{Solution 3 (substitution)}

Let $1-x=a;1-y=b;1-z=c$, rewrite those equations

 $\sqrt{(1-a)(1+b)}+\sqrt{(1+a)(1-b)}=1$;

 $\sqrt{(1-b)(1+c)}+\sqrt{(1+b)(1-c)}=\sqrt{2}$

 $\sqrt{(1-a)(1+c)}+\sqrt{(1-c)(1+a)}=\sqrt{3}$

 and solve for $m/n = (abc)^2 = a^2b^2c^2$

 Square both sides and simplify, to get three equations:

 $2ab-1=2\sqrt{(1-a^2)(1-b^2)}$

 $2bc~ ~ ~ ~  ~ ~=2\sqrt{(1-b^2)(1-c^2)}$

 $2ac+1=2\sqrt{(1-c^2)(1-a^2)}$

 Square both sides again, and simplify to get three equations:

 $a^2+b^2-ab=\frac{3}{4}$

 $b^2+c^2~ ~ ~ ~ ~ ~=1$

 $a^2+c^2+ac=\frac{3}{4}$

 Subtract first and third equation, getting $(b+c)(b-c)=a(b+c)$, $a=b-c$

 Put it in first equation, getting $b^2-2bc+c^2+b^2-b(b-c)=b^2+c^2-bc=\frac{3}{4}$, $bc=\frac{1}{4}$

 Since $a^2=b^2+c^2-2bc=\frac{1}{2}$, $m/n = a^2b^2c^2 = a^2(bc)^2 = \frac{1}{2}\left(\frac{1}{4}\right)^2=\frac{1}{32}$  and so the final answer is $\boxed{033}$

 





\textbf{Solution 4}

Denote $u = 1 - x$, $v = 1 - y$, $w = 1 - z$.Hence, the system of equations given in the problem can be written as
 \begin{align*} \sqrt{(1-u)(1+v)} + \sqrt{(1+u)(1-v)} & = 1 \hspace{1cm} (1) \\ \sqrt{(1-v)(1+w)} + \sqrt{(1+v)(1-w)} & = \sqrt{2} \hspace{1cm} (2) \\ \sqrt{(1-w)(1+u)} + \sqrt{(1+w)(1-u)} & = \sqrt{3} . \hspace{1cm} (3)  \end{align*}


 Each equation above takes the following form:
 \[ \sqrt{(1-a)(1+b)} + \sqrt{(1+a)(1-b)} = k . \]

 Now, we simplify this equation by removing radicals.

 Denote $p = \sqrt{(1-a)(1+b)}$ and $q = \sqrt{(1+a)(1-b)}$.

 Hence, the equation above implies
 \[ \left\{ \begin{array}{l} p + q = k \\ p^2 = (1-a)(1+b) \\ q^2 = (1+a)(1-b) \end{array} \right.. \]

 Hence, $q^2 - p^2 = (1+a)(1-b) - (1-a)(1+b) = 2 (a-b)$.Hence, $q - p = \frac{q^2 - p^2}{p+q} = \frac{2}{k} (a-b)$.

 Because $p + q = k$ and $q - p = \frac{2}{k} (a-b)$, we get $q = \frac{a-b}{k} + \frac{k}{2}$.Plugging this into the equation $q^2 = (1+a)(1-b)$ and simplifying it, we get
 \[ a^2 + \left( k^2 - 2 \right) ab + b^2 = k^2 - \frac{k^4}{4} . \]

 Therefore, the system of equations above can be simplified as
 \begin{align*} u^2 - uv + v^2 & = \frac{3}{4} \\ v^2 + w^2 & = 1 \\ w^2 + wu + u^2 & = \frac{3}{4}  . \end{align*}


 Denote $w' = - w$.The system of equations above can be equivalently written as
 \begin{align*} u^2 - uv + v^2 & = \frac{3}{4} \hspace{1cm} (1') \\ v^2 + w'^2 & = 1 \hspace{1cm} (2') \\ w'^2 - w'u + u^2 & = \frac{3}{4} \hspace{1cm} (3') . \end{align*}


 Taking $(1') - (3')$, we get
 \[ (v - w') (v + w' - u) = 0 . \]

 Thus, we have either $v - w' = 0$ or $v + w' - u = 0$.

 $\textbf{Case 1}$: $v - w' = 0$.

 Equation (2') implies $v = w' = \pm \frac{1}{\sqrt{2}}$.

 Plugging $v$ and $w'$ into Equation (2), we get contradiction. Therefore, this case is infeasible.

 $\textbf{Case 2}$: $v + w' - u = 0$.

 Plugging this condition into (1') to substitute $u$, we get
 \[ v^2 + v w' + w'^2 = \frac{3}{4} \hspace{1cm} (4) . \]

 Taking $(4) - (2')$, we get
 \[ v w' = - \frac{1}{4} . \hspace{1cm} (5) . \]

 Taking (4) + (5), we get
 \[ \left( v + w' \right)^2 = \frac{1}{2} . \]

 Hence, $u^2 = \left( v + w' \right)^2 = \frac{1}{2}$.

 Therefore,
 \begin{align*} \left[ (1-x)(1-y)(1-z) \right]^2 & = u^2 (vw)^2 \\ & = u^2 (vw')^2 \\ & = \frac{1}{2} \left( - \frac{1}{4} \right)^2 \\ & = \frac{1}{32} . \end{align*}


 Therefore, the answer is $1 + 32 = \boxed{\textbf{(033) }}$.

 bu-bye
 \begin{align*} \sin(\alpha + \beta) &= 1/2 \\  \sin(\alpha + \gamma) &= \sqrt2/2 \\  \sin(\beta + \gamma) &= \sqrt3/2. \end{align*}
Thus, 
 \begin{align*} \alpha + \beta &= 30^{\circ} \\  \alpha + \gamma &= 45^{\circ} \\  \beta + \gamma &= 60^{\circ}, \end{align*}
so $(\alpha, \beta, \gamma) = (15/2^{\circ}, 45/2^{\circ}, 75/2^{\circ})$. Hence,

 
 \[abc = (1-2\sin^2(\alpha))(1-2\sin^2(\beta))(1-2\sin^2(\gamma))=\cos(15^{\circ})\cos(45^{\circ})\cos(75^{\circ})=\frac{\sqrt{2}}{8},\]so $(abc)^2=(\sqrt{2}/8)^2=\frac{1}{32}$, for a final answer of $\boxed{033}$.

 <u>\textbf{Remark}</u>

 The motivation for the trig substitution is that if $\sin^2(\alpha)=(1-a)/2$, then $\cos^2(\alpha)=(1+a)/2$, and when making the substitution in each equation of the initial set of equations, we obtain a new equation in the form of the sine addition formula.

 





\textbf{Solution 6 (Geometric)}


                    \begin{figure}[H]
                        \centering
                        \includegraphics[width=0.5\textwidth]{Solution/AIME_2022_I_15_Solution_Figure_1.png}
                    \end{figure}
                    In given equations, $0 \leq x,y,z \leq 2,$ so we define some points:
 \[\bar {O} = (0, 0), \bar {A} = (1, 0), \bar{M} = \left(\frac {1}{\sqrt{2}},\frac {1}{\sqrt{2}}\right),\] 
 \[\bar {X} = \left(\sqrt {\frac {x}{2}}, \sqrt{1 – \frac{x}{2}}\right), \bar {Y'} = \left(\sqrt {\frac {y}{2}}, \sqrt{1 – \frac{y}{2}}\right),\]
 \[\bar {Y} = \left(\sqrt {1 – \frac{y}{2}},\sqrt{\frac {y}{2}}\right), \bar {Z} = \left(\sqrt {1 – \frac{z}{2}},\sqrt{\frac {z}{2}}\right).\]Notice, that 
 \[\mid \vec {AO} \mid = \mid \vec {MO} \mid = \mid \vec {XO} \mid =\mid \vec {YO} \mid = \mid \vec {Y'O} \mid =\mid \vec {ZO} \mid = 1\] and each points lies in the first quadrant.

 We use given equations and get some scalar products:
 \[(\vec {XO} \cdot \vec {YO}) = \frac {1}{2} = \cos \angle XOY \implies  \angle XOY = 60 ^\circ,\] 
 \[(\vec {XO} \cdot \vec {ZO}) = \frac {\sqrt{3}}{2} = \cos \angle XOZ \implies  \angle XOZ = 30^\circ,\] 
 \[(\vec {Y'O} \cdot \vec {ZO}) = \frac {1}{\sqrt{2}} = \cos \angle Y'OZ \implies  \angle Y'OZ = 45^\circ.\] So $\angle YOZ =  \angle XOY – \angle XOZ =  60 ^\circ – 30 ^\circ = 30 ^\circ,   \angle Y'OY =  \angle Y'OZ +  \angle YOZ = 45^\circ + 30 ^\circ = 75^\circ.$

 Points $Y$ and $Y'$ are symmetric with respect to $OM.$

 <i>\textbf{Case 1}</i>
 \[\angle YOA = \frac{90^\circ – 75^\circ}{2} = 7.5^\circ, \angle ZOA = 30^\circ + 7.5^\circ = 37.5^\circ, \angle XOA = 60^\circ + 7.5^\circ = 67.5^\circ .\]
 \[1 – x = \left(\sqrt{1 – \frac{x}{2}} \right)^2– \left(\sqrt{\frac {x}{2}}\right)^2  = \sin^2 \angle XOA – \cos^2 \angle XOA = –\cos 2 \angle XOA = –\cos 135^\circ,\]
 \[1 – y = \cos 15^\circ, 1 – z = \cos 75^\circ \implies  \left[ (1–x)(1–y)(1–z) \right]^2 = \left[ \sin 45^\circ \cdot \cos 15^\circ \cdot \sin 15^\circ \right]^2 =\]
 \[=\left[ \frac {\sin 45^\circ \cdot \sin 30^\circ}{2} \right]^2  = \frac {1}{32} \implies \boxed{\textbf{033}}.\]<i>\textbf{Case 2}</i>

 
 \[\angle Y_1 OA = \frac{90^\circ + 75^\circ}{2} = 82.5^\circ, \angle Z_1 OA = 82.5^\circ – 30^\circ = 52.5^\circ, \angle X_1 OA = 82.5^\circ – 60^\circ = 22.5^\circ \implies \boxed{\textbf{033}}.\]

 \textbf{vladimir.shelomovskii@gmail.com, vvsss}

 


 


 




\end{enumerate}